\documentclass[11pt,a4paper]{article}
% ============================================
% GIFT v3.4 — Shared Preamble
% ============================================
% Usage: % ============================================
% GIFT v3.4 — Shared Preamble
% ============================================
% Usage: % ============================================
% GIFT v3.4 — Shared Preamble
% ============================================
% Usage: \input{gift_preamble} after \documentclass

% ENCODING & FONTS
\usepackage[utf8]{inputenc}
\usepackage[T1]{fontenc}
\usepackage{lmodern}

% PAGE LAYOUT (Golden Ratio)
\usepackage[margin=1.618cm, top=2.618cm, bottom=2.618cm]{geometry}

% ESSENTIAL PACKAGES
\usepackage{float}
\usepackage{caption}
\usepackage{subcaption}
\usepackage{setspace}
\usepackage{fancyhdr}
\usepackage{xcolor}
\usepackage{hyperref}
\usepackage{csquotes}
\usepackage{amsmath}
\usepackage{amssymb}
\usepackage{amsthm}
\usepackage{mathtools}
\usepackage{booktabs}
\usepackage{longtable}
\usepackage{array}
\usepackage{multirow}
\usepackage{tikz}
\usepackage{graphicx}
\usepackage{listings}
\usepackage{enumitem}
% Unicode fallbacks (pandoc artifacts that survive markdown→LaTeX)
\DeclareUnicodeCharacter{00B0}{\ensuremath{^\circ}}
\DeclareUnicodeCharacter{221A}{\ensuremath{\sqrt{\,}}}
\DeclareUnicodeCharacter{2713}{\ensuremath{\checkmark}}
\DeclareUnicodeCharacter{2717}{\ensuremath{\times}}
\DeclareUnicodeCharacter{2070}{\ensuremath{^{0}}}
\DeclareUnicodeCharacter{00B9}{\ensuremath{^{1}}}
\DeclareUnicodeCharacter{00B2}{\ensuremath{^{2}}}
\DeclareUnicodeCharacter{00B3}{\ensuremath{^{3}}}
\DeclareUnicodeCharacter{2074}{\ensuremath{^{4}}}
\DeclareUnicodeCharacter{2075}{\ensuremath{^{5}}}
\DeclareUnicodeCharacter{2076}{\ensuremath{^{6}}}
\DeclareUnicodeCharacter{2077}{\ensuremath{^{7}}}
\DeclareUnicodeCharacter{2078}{\ensuremath{^{8}}}
\DeclareUnicodeCharacter{2079}{\ensuremath{^{9}}}
\DeclareUnicodeCharacter{2071}{\ensuremath{^{i}}}
\DeclareUnicodeCharacter{2080}{\ensuremath{_{0}}}
\DeclareUnicodeCharacter{2081}{\ensuremath{_{1}}}
\DeclareUnicodeCharacter{2082}{\ensuremath{_{2}}}
\DeclareUnicodeCharacter{2083}{\ensuremath{_{3}}}
\DeclareUnicodeCharacter{2084}{\ensuremath{_{4}}}
\DeclareUnicodeCharacter{2085}{\ensuremath{_{5}}}
\DeclareUnicodeCharacter{2086}{\ensuremath{_{6}}}
\DeclareUnicodeCharacter{2087}{\ensuremath{_{7}}}
\DeclareUnicodeCharacter{2088}{\ensuremath{_{8}}}
\DeclareUnicodeCharacter{2089}{\ensuremath{_{9}}}
\DeclareUnicodeCharacter{208A}{\ensuremath{_{+}}}
\DeclareUnicodeCharacter{208B}{\ensuremath{_{-}}}
\DeclareUnicodeCharacter{2192}{\ensuremath{\to}}
\DeclareUnicodeCharacter{2502}{|}
\DeclareUnicodeCharacter{25BC}{\ensuremath{\blacktriangledown}}
\DeclareUnicodeCharacter{2500}{\ensuremath{-}}
\DeclareUnicodeCharacter{2514}{\ensuremath{\llcorner}}
\DeclareUnicodeCharacter{251C}{\ensuremath{\vdash}}
\DeclareUnicodeCharacter{2026}{\ldots}
\DeclareUnicodeCharacter{2208}{\ensuremath{\in}}
\DeclareUnicodeCharacter{2229}{\ensuremath{\cap}}
\DeclareUnicodeCharacter{222A}{\ensuremath{\cup}}
\DeclareUnicodeCharacter{2295}{\ensuremath{\oplus}}
\DeclareUnicodeCharacter{2297}{\ensuremath{\otimes}}
\DeclareUnicodeCharacter{2245}{\ensuremath{\cong}}
\DeclareUnicodeCharacter{2248}{\ensuremath{\approx}}
\DeclareUnicodeCharacter{2260}{\ensuremath{\neq}}
\DeclareUnicodeCharacter{2261}{\ensuremath{\equiv}}
\DeclareUnicodeCharacter{2264}{\ensuremath{\leq}}
\DeclareUnicodeCharacter{2265}{\ensuremath{\geq}}
\DeclareUnicodeCharacter{00D7}{\ensuremath{\times}}
\DeclareUnicodeCharacter{00B1}{\ensuremath{\pm}}
\DeclareUnicodeCharacter{2225}{\ensuremath{\|}}
% Greek letters used in body text (pandoc passes them through unchanged)
\DeclareUnicodeCharacter{03B1}{\ensuremath{\alpha}}
\DeclareUnicodeCharacter{03B2}{\ensuremath{\beta}}
\DeclareUnicodeCharacter{03B3}{\ensuremath{\gamma}}
\DeclareUnicodeCharacter{03B4}{\ensuremath{\delta}}
\DeclareUnicodeCharacter{03B5}{\ensuremath{\varepsilon}}
\DeclareUnicodeCharacter{03B6}{\ensuremath{\zeta}}
\DeclareUnicodeCharacter{03B7}{\ensuremath{\eta}}
\DeclareUnicodeCharacter{03B8}{\ensuremath{\theta}}
\DeclareUnicodeCharacter{03B9}{\ensuremath{\iota}}
\DeclareUnicodeCharacter{03BA}{\ensuremath{\kappa}}
\DeclareUnicodeCharacter{03BB}{\ensuremath{\lambda}}
\DeclareUnicodeCharacter{03BC}{\ensuremath{\mu}}
\DeclareUnicodeCharacter{03BD}{\ensuremath{\nu}}
\DeclareUnicodeCharacter{03BE}{\ensuremath{\xi}}
\DeclareUnicodeCharacter{03C0}{\ensuremath{\pi}}
\DeclareUnicodeCharacter{03C1}{\ensuremath{\rho}}
\DeclareUnicodeCharacter{03C3}{\ensuremath{\sigma}}
\DeclareUnicodeCharacter{03C4}{\ensuremath{\tau}}
\DeclareUnicodeCharacter{03C5}{\ensuremath{\upsilon}}
\DeclareUnicodeCharacter{03C6}{\ensuremath{\varphi}}
\DeclareUnicodeCharacter{03C7}{\ensuremath{\chi}}
\DeclareUnicodeCharacter{03C8}{\ensuremath{\psi}}
\DeclareUnicodeCharacter{03C9}{\ensuremath{\omega}}
\DeclareUnicodeCharacter{0394}{\ensuremath{\Delta}}
\DeclareUnicodeCharacter{0398}{\ensuremath{\Theta}}
\DeclareUnicodeCharacter{039B}{\ensuremath{\Lambda}}
\DeclareUnicodeCharacter{03A0}{\ensuremath{\Pi}}
\DeclareUnicodeCharacter{03A3}{\ensuremath{\Sigma}}
\DeclareUnicodeCharacter{03A6}{\ensuremath{\Phi}}
\DeclareUnicodeCharacter{03A8}{\ensuremath{\Psi}}
\DeclareUnicodeCharacter{03A9}{\ensuremath{\Omega}}
\DeclareUnicodeCharacter{2200}{\ensuremath{\forall}}
\DeclareUnicodeCharacter{2203}{\ensuremath{\exists}}
\DeclareUnicodeCharacter{2207}{\ensuremath{\nabla}}
\DeclareUnicodeCharacter{2202}{\ensuremath{\partial}}
\DeclareUnicodeCharacter{2218}{\ensuremath{\circ}}
\DeclareUnicodeCharacter{2032}{\ensuremath{\prime}}
\DeclareUnicodeCharacter{2113}{\ensuremath{\ell}}
\DeclareUnicodeCharacter{210F}{\ensuremath{\hbar}}
\DeclareUnicodeCharacter{27E8}{\ensuremath{\langle}}
\DeclareUnicodeCharacter{27E9}{\ensuremath{\rangle}}
\DeclareUnicodeCharacter{2308}{\ensuremath{\lceil}}
\DeclareUnicodeCharacter{2309}{\ensuremath{\rceil}}
\DeclareUnicodeCharacter{230A}{\ensuremath{\lfloor}}
\DeclareUnicodeCharacter{230B}{\ensuremath{\rfloor}}
\DeclareUnicodeCharacter{222B}{\ensuremath{\int}}
\DeclareUnicodeCharacter{2211}{\ensuremath{\sum}}
\DeclareUnicodeCharacter{220F}{\ensuremath{\prod}}
\DeclareUnicodeCharacter{221E}{\ensuremath{\infty}}
\DeclareUnicodeCharacter{00B5}{\ensuremath{\mu}}
\DeclareUnicodeCharacter{2236}{\ensuremath{\colon}}
\DeclareUnicodeCharacter{2192}{\ensuremath{\to}}
\DeclareUnicodeCharacter{21A6}{\ensuremath{\mapsto}}
\DeclareUnicodeCharacter{21D2}{\ensuremath{\Rightarrow}}
\DeclareUnicodeCharacter{21D4}{\ensuremath{\Leftrightarrow}}
\DeclareUnicodeCharacter{2282}{\ensuremath{\subset}}
\DeclareUnicodeCharacter{2286}{\ensuremath{\subseteq}}
\DeclareUnicodeCharacter{2287}{\ensuremath{\supseteq}}
\DeclareUnicodeCharacter{222B}{\ensuremath{\int}}
\DeclareUnicodeCharacter{2205}{\ensuremath{\emptyset}}
\DeclareUnicodeCharacter{2299}{\ensuremath{\odot}}
\DeclareUnicodeCharacter{2220}{\ensuremath{\angle}}
\DeclareUnicodeCharacter{2135}{\ensuremath{\aleph}}
\DeclareUnicodeCharacter{207B}{\ensuremath{^{-}}}
\DeclareUnicodeCharacter{207A}{\ensuremath{^{+}}}
\DeclareUnicodeCharacter{220E}{\ensuremath{\blacksquare}}
\DeclareUnicodeCharacter{2061}{}
\DeclareUnicodeCharacter{2062}{}
\DeclareUnicodeCharacter{2063}{}
\DeclareUnicodeCharacter{2009}{\,}
\DeclareUnicodeCharacter{200A}{\,}
\DeclareUnicodeCharacter{200B}{}
\DeclareUnicodeCharacter{2212}{\ensuremath{-}}
\DeclareUnicodeCharacter{2194}{\ensuremath{\leftrightarrow}}
\DeclareUnicodeCharacter{2099}{\ensuremath{_{n}}}
\DeclareUnicodeCharacter{2098}{\ensuremath{_{m}}}
\DeclareUnicodeCharacter{2096}{\ensuremath{_{k}}}
\DeclareUnicodeCharacter{2095}{\ensuremath{_{h}}}
\DeclareUnicodeCharacter{2090}{\ensuremath{_{a}}}
\DeclareUnicodeCharacter{2091}{\ensuremath{_{e}}}
\DeclareUnicodeCharacter{2093}{\ensuremath{_{x}}}
\DeclareUnicodeCharacter{1D62}{\ensuremath{_{i}}}
\DeclareUnicodeCharacter{2C7C}{\ensuremath{_{j}}}
\DeclareUnicodeCharacter{2C7B}{\ensuremath{^{e}}}
\DeclareUnicodeCharacter{1D52}{\ensuremath{^{o}}}
\DeclareUnicodeCharacter{1D5}{\ensuremath{^{u}}}
\DeclareUnicodeCharacter{02E1}{\ensuremath{^{l}}}
\DeclareUnicodeCharacter{2295}{\ensuremath{\oplus}}
\DeclareUnicodeCharacter{27FA}{\ensuremath{\Longleftrightarrow}}
\DeclareUnicodeCharacter{27F9}{\ensuremath{\Longrightarrow}}
\DeclareUnicodeCharacter{2299}{\ensuremath{\odot}}
% Provide \tightlist used by pandoc-generated lists (no-op)
\providecommand{\tightlist}{\setlength{\itemsep}{0pt}\setlength{\parskip}{0pt}}
% pandoc longtable column specs: provide \real as passthrough and a
% `none` counter so that calc's \real{x} expansion path works.
\newcounter{none}
\providecommand{\real}[1]{#1}

% LISTINGS
\lstset{
    basicstyle=\small\ttfamily,
    breaklines=true, frame=single,
    keepspaces=true, showstringspaces=false,
    aboveskip=0.8em, belowskip=0.8em
}

% HEADER/FOOTER
\setlength{\headheight}{14pt}
\pagestyle{fancy}
\fancyhf{}
\fancyhead[R]{\thepage}
\renewcommand{\headrulewidth}{0.2pt}

% HYPERREF
\hypersetup{
    colorlinks=true,
    linkcolor=blue, citecolor=blue, urlcolor=blue,
    pdfauthor={Brieuc de La Fourniere}
}

% SPACING
\setstretch{1.2}
\setlength{\parskip}{0.4em}
\setlength{\parindent}{0pt}

% TITLE FORMATTING
\usepackage{titling}
\pretitle{\LARGE\bfseries}
\posttitle{\vspace{-0.4em}}
\preauthor{}
\postauthor{}
\predate{}
\postdate{}
\setlength{\droptitle}{-2.0em}

% CUSTOM COMMANDS — GIFT notation
\newcommand{\E}{\mathrm{E}}
\newcommand{\Gtwo}{\mathrm{G}_2}
\newcommand{\Kseven}{K_7}
\newcommand{\AdS}{\mathrm{AdS}}
\newcommand{\SU}{\mathrm{SU}}
\newcommand{\SO}{\mathrm{SO}}
\newcommand{\U}{\mathrm{U}}
\newcommand{\Sp}{\mathrm{Sp}}
\newcommand{\PSL}{\mathrm{PSL}}
\newcommand{\SM}{\mathrm{SM}}
\newcommand{\Tr}{\mathrm{Tr}}
\newcommand{\rank}{\mathrm{rank}}
\newcommand{\Vol}{\mathrm{Vol}}
\newcommand{\Hol}{\mathrm{Hol}}
\newcommand{\Ric}{\mathrm{Ric}}
\newcommand{\Rm}{\mathrm{Rm}}
\newcommand{\Scal}{\mathrm{Scal}}
\DeclareMathOperator{\diag}{diag}
\DeclareMathOperator{\cond}{cond}
\newcommand{\GIFT}{\textrm{GIFT}}
\newcommand{\CP}{\mathrm{CP}}
\newcommand{\CKM}{\mathrm{CKM}}
\newcommand{\PMNS}{\mathrm{PMNS}}
\newcommand{\EW}{\mathrm{EW}}
\newcommand{\Pl}{\mathrm{Pl}}
\newcommand{\DE}{\mathrm{DE}}
\newcommand{\DM}{\mathrm{DM}}
\newcommand{\KK}{\mathrm{KK}}
\newcommand{\NK}{\mathrm{NK}}
\newcommand{\TCS}{\mathrm{TCS}}
\newcommand{\PINN}{\mathrm{PINN}}
\newcommand{\btwo}{b_2}
\newcommand{\bthree}{b_3}
\newcommand{\typeI}{\textsc{I}}
\newcommand{\typeII}{\textsc{II}}
\newcommand{\typeIII}{\textsc{III}}
\newcommand{\typeIV}{\textsc{IV}}
\newcommand{\proven}{\textsc{Proven}}

% PDF bookmark safety
\pdfstringdefDisableCommands{%
  \def\textsubscript#1{#1}%
  \def\textsuperscript#1{#1}%
  \def\CP{CP}%
  \def\Gtwo{G2}%
  \def\Kseven{K7}%
  \def\GIFT{GIFT}%
  \def\E{E}%
  \def\SM{SM}%
}

% THEOREM ENVIRONMENTS
\newtheorem{theorem}{Theorem}[section]
\newtheorem{proposition}[theorem]{Proposition}
\newtheorem{lemma}[theorem]{Lemma}
\newtheorem{corollary}[theorem]{Corollary}
\theoremstyle{definition}
\newtheorem{definition}[theorem]{Definition}
\theoremstyle{remark}
\newtheorem{remark}[theorem]{Remark}
 after \documentclass

% ENCODING & FONTS
\usepackage[utf8]{inputenc}
\usepackage[T1]{fontenc}
\usepackage{lmodern}

% PAGE LAYOUT (Golden Ratio)
\usepackage[margin=1.618cm, top=2.618cm, bottom=2.618cm]{geometry}

% ESSENTIAL PACKAGES
\usepackage{float}
\usepackage{caption}
\usepackage{subcaption}
\usepackage{setspace}
\usepackage{fancyhdr}
\usepackage{xcolor}
\usepackage{hyperref}
\usepackage{csquotes}
\usepackage{amsmath}
\usepackage{amssymb}
\usepackage{amsthm}
\usepackage{mathtools}
\usepackage{booktabs}
\usepackage{longtable}
\usepackage{array}
\usepackage{multirow}
\usepackage{tikz}
\usepackage{graphicx}
\usepackage{listings}
\usepackage{enumitem}
% Unicode fallbacks (pandoc artifacts that survive markdown→LaTeX)
\DeclareUnicodeCharacter{00B0}{\ensuremath{^\circ}}
\DeclareUnicodeCharacter{221A}{\ensuremath{\sqrt{\,}}}
\DeclareUnicodeCharacter{2713}{\ensuremath{\checkmark}}
\DeclareUnicodeCharacter{2717}{\ensuremath{\times}}
\DeclareUnicodeCharacter{2070}{\ensuremath{^{0}}}
\DeclareUnicodeCharacter{00B9}{\ensuremath{^{1}}}
\DeclareUnicodeCharacter{00B2}{\ensuremath{^{2}}}
\DeclareUnicodeCharacter{00B3}{\ensuremath{^{3}}}
\DeclareUnicodeCharacter{2074}{\ensuremath{^{4}}}
\DeclareUnicodeCharacter{2075}{\ensuremath{^{5}}}
\DeclareUnicodeCharacter{2076}{\ensuremath{^{6}}}
\DeclareUnicodeCharacter{2077}{\ensuremath{^{7}}}
\DeclareUnicodeCharacter{2078}{\ensuremath{^{8}}}
\DeclareUnicodeCharacter{2079}{\ensuremath{^{9}}}
\DeclareUnicodeCharacter{2071}{\ensuremath{^{i}}}
\DeclareUnicodeCharacter{2080}{\ensuremath{_{0}}}
\DeclareUnicodeCharacter{2081}{\ensuremath{_{1}}}
\DeclareUnicodeCharacter{2082}{\ensuremath{_{2}}}
\DeclareUnicodeCharacter{2083}{\ensuremath{_{3}}}
\DeclareUnicodeCharacter{2084}{\ensuremath{_{4}}}
\DeclareUnicodeCharacter{2085}{\ensuremath{_{5}}}
\DeclareUnicodeCharacter{2086}{\ensuremath{_{6}}}
\DeclareUnicodeCharacter{2087}{\ensuremath{_{7}}}
\DeclareUnicodeCharacter{2088}{\ensuremath{_{8}}}
\DeclareUnicodeCharacter{2089}{\ensuremath{_{9}}}
\DeclareUnicodeCharacter{208A}{\ensuremath{_{+}}}
\DeclareUnicodeCharacter{208B}{\ensuremath{_{-}}}
\DeclareUnicodeCharacter{2192}{\ensuremath{\to}}
\DeclareUnicodeCharacter{2502}{|}
\DeclareUnicodeCharacter{25BC}{\ensuremath{\blacktriangledown}}
\DeclareUnicodeCharacter{2500}{\ensuremath{-}}
\DeclareUnicodeCharacter{2514}{\ensuremath{\llcorner}}
\DeclareUnicodeCharacter{251C}{\ensuremath{\vdash}}
\DeclareUnicodeCharacter{2026}{\ldots}
\DeclareUnicodeCharacter{2208}{\ensuremath{\in}}
\DeclareUnicodeCharacter{2229}{\ensuremath{\cap}}
\DeclareUnicodeCharacter{222A}{\ensuremath{\cup}}
\DeclareUnicodeCharacter{2295}{\ensuremath{\oplus}}
\DeclareUnicodeCharacter{2297}{\ensuremath{\otimes}}
\DeclareUnicodeCharacter{2245}{\ensuremath{\cong}}
\DeclareUnicodeCharacter{2248}{\ensuremath{\approx}}
\DeclareUnicodeCharacter{2260}{\ensuremath{\neq}}
\DeclareUnicodeCharacter{2261}{\ensuremath{\equiv}}
\DeclareUnicodeCharacter{2264}{\ensuremath{\leq}}
\DeclareUnicodeCharacter{2265}{\ensuremath{\geq}}
\DeclareUnicodeCharacter{00D7}{\ensuremath{\times}}
\DeclareUnicodeCharacter{00B1}{\ensuremath{\pm}}
\DeclareUnicodeCharacter{2225}{\ensuremath{\|}}
% Greek letters used in body text (pandoc passes them through unchanged)
\DeclareUnicodeCharacter{03B1}{\ensuremath{\alpha}}
\DeclareUnicodeCharacter{03B2}{\ensuremath{\beta}}
\DeclareUnicodeCharacter{03B3}{\ensuremath{\gamma}}
\DeclareUnicodeCharacter{03B4}{\ensuremath{\delta}}
\DeclareUnicodeCharacter{03B5}{\ensuremath{\varepsilon}}
\DeclareUnicodeCharacter{03B6}{\ensuremath{\zeta}}
\DeclareUnicodeCharacter{03B7}{\ensuremath{\eta}}
\DeclareUnicodeCharacter{03B8}{\ensuremath{\theta}}
\DeclareUnicodeCharacter{03B9}{\ensuremath{\iota}}
\DeclareUnicodeCharacter{03BA}{\ensuremath{\kappa}}
\DeclareUnicodeCharacter{03BB}{\ensuremath{\lambda}}
\DeclareUnicodeCharacter{03BC}{\ensuremath{\mu}}
\DeclareUnicodeCharacter{03BD}{\ensuremath{\nu}}
\DeclareUnicodeCharacter{03BE}{\ensuremath{\xi}}
\DeclareUnicodeCharacter{03C0}{\ensuremath{\pi}}
\DeclareUnicodeCharacter{03C1}{\ensuremath{\rho}}
\DeclareUnicodeCharacter{03C3}{\ensuremath{\sigma}}
\DeclareUnicodeCharacter{03C4}{\ensuremath{\tau}}
\DeclareUnicodeCharacter{03C5}{\ensuremath{\upsilon}}
\DeclareUnicodeCharacter{03C6}{\ensuremath{\varphi}}
\DeclareUnicodeCharacter{03C7}{\ensuremath{\chi}}
\DeclareUnicodeCharacter{03C8}{\ensuremath{\psi}}
\DeclareUnicodeCharacter{03C9}{\ensuremath{\omega}}
\DeclareUnicodeCharacter{0394}{\ensuremath{\Delta}}
\DeclareUnicodeCharacter{0398}{\ensuremath{\Theta}}
\DeclareUnicodeCharacter{039B}{\ensuremath{\Lambda}}
\DeclareUnicodeCharacter{03A0}{\ensuremath{\Pi}}
\DeclareUnicodeCharacter{03A3}{\ensuremath{\Sigma}}
\DeclareUnicodeCharacter{03A6}{\ensuremath{\Phi}}
\DeclareUnicodeCharacter{03A8}{\ensuremath{\Psi}}
\DeclareUnicodeCharacter{03A9}{\ensuremath{\Omega}}
\DeclareUnicodeCharacter{2200}{\ensuremath{\forall}}
\DeclareUnicodeCharacter{2203}{\ensuremath{\exists}}
\DeclareUnicodeCharacter{2207}{\ensuremath{\nabla}}
\DeclareUnicodeCharacter{2202}{\ensuremath{\partial}}
\DeclareUnicodeCharacter{2218}{\ensuremath{\circ}}
\DeclareUnicodeCharacter{2032}{\ensuremath{\prime}}
\DeclareUnicodeCharacter{2113}{\ensuremath{\ell}}
\DeclareUnicodeCharacter{210F}{\ensuremath{\hbar}}
\DeclareUnicodeCharacter{27E8}{\ensuremath{\langle}}
\DeclareUnicodeCharacter{27E9}{\ensuremath{\rangle}}
\DeclareUnicodeCharacter{2308}{\ensuremath{\lceil}}
\DeclareUnicodeCharacter{2309}{\ensuremath{\rceil}}
\DeclareUnicodeCharacter{230A}{\ensuremath{\lfloor}}
\DeclareUnicodeCharacter{230B}{\ensuremath{\rfloor}}
\DeclareUnicodeCharacter{222B}{\ensuremath{\int}}
\DeclareUnicodeCharacter{2211}{\ensuremath{\sum}}
\DeclareUnicodeCharacter{220F}{\ensuremath{\prod}}
\DeclareUnicodeCharacter{221E}{\ensuremath{\infty}}
\DeclareUnicodeCharacter{00B5}{\ensuremath{\mu}}
\DeclareUnicodeCharacter{2236}{\ensuremath{\colon}}
\DeclareUnicodeCharacter{2192}{\ensuremath{\to}}
\DeclareUnicodeCharacter{21A6}{\ensuremath{\mapsto}}
\DeclareUnicodeCharacter{21D2}{\ensuremath{\Rightarrow}}
\DeclareUnicodeCharacter{21D4}{\ensuremath{\Leftrightarrow}}
\DeclareUnicodeCharacter{2282}{\ensuremath{\subset}}
\DeclareUnicodeCharacter{2286}{\ensuremath{\subseteq}}
\DeclareUnicodeCharacter{2287}{\ensuremath{\supseteq}}
\DeclareUnicodeCharacter{222B}{\ensuremath{\int}}
\DeclareUnicodeCharacter{2205}{\ensuremath{\emptyset}}
\DeclareUnicodeCharacter{2299}{\ensuremath{\odot}}
\DeclareUnicodeCharacter{2220}{\ensuremath{\angle}}
\DeclareUnicodeCharacter{2135}{\ensuremath{\aleph}}
\DeclareUnicodeCharacter{207B}{\ensuremath{^{-}}}
\DeclareUnicodeCharacter{207A}{\ensuremath{^{+}}}
\DeclareUnicodeCharacter{220E}{\ensuremath{\blacksquare}}
\DeclareUnicodeCharacter{2061}{}
\DeclareUnicodeCharacter{2062}{}
\DeclareUnicodeCharacter{2063}{}
\DeclareUnicodeCharacter{2009}{\,}
\DeclareUnicodeCharacter{200A}{\,}
\DeclareUnicodeCharacter{200B}{}
\DeclareUnicodeCharacter{2212}{\ensuremath{-}}
\DeclareUnicodeCharacter{2194}{\ensuremath{\leftrightarrow}}
\DeclareUnicodeCharacter{2099}{\ensuremath{_{n}}}
\DeclareUnicodeCharacter{2098}{\ensuremath{_{m}}}
\DeclareUnicodeCharacter{2096}{\ensuremath{_{k}}}
\DeclareUnicodeCharacter{2095}{\ensuremath{_{h}}}
\DeclareUnicodeCharacter{2090}{\ensuremath{_{a}}}
\DeclareUnicodeCharacter{2091}{\ensuremath{_{e}}}
\DeclareUnicodeCharacter{2093}{\ensuremath{_{x}}}
\DeclareUnicodeCharacter{1D62}{\ensuremath{_{i}}}
\DeclareUnicodeCharacter{2C7C}{\ensuremath{_{j}}}
\DeclareUnicodeCharacter{2C7B}{\ensuremath{^{e}}}
\DeclareUnicodeCharacter{1D52}{\ensuremath{^{o}}}
\DeclareUnicodeCharacter{1D5}{\ensuremath{^{u}}}
\DeclareUnicodeCharacter{02E1}{\ensuremath{^{l}}}
\DeclareUnicodeCharacter{2295}{\ensuremath{\oplus}}
\DeclareUnicodeCharacter{27FA}{\ensuremath{\Longleftrightarrow}}
\DeclareUnicodeCharacter{27F9}{\ensuremath{\Longrightarrow}}
\DeclareUnicodeCharacter{2299}{\ensuremath{\odot}}
% Provide \tightlist used by pandoc-generated lists (no-op)
\providecommand{\tightlist}{\setlength{\itemsep}{0pt}\setlength{\parskip}{0pt}}
% pandoc longtable column specs: provide \real as passthrough and a
% `none` counter so that calc's \real{x} expansion path works.
\newcounter{none}
\providecommand{\real}[1]{#1}

% LISTINGS
\lstset{
    basicstyle=\small\ttfamily,
    breaklines=true, frame=single,
    keepspaces=true, showstringspaces=false,
    aboveskip=0.8em, belowskip=0.8em
}

% HEADER/FOOTER
\setlength{\headheight}{14pt}
\pagestyle{fancy}
\fancyhf{}
\fancyhead[R]{\thepage}
\renewcommand{\headrulewidth}{0.2pt}

% HYPERREF
\hypersetup{
    colorlinks=true,
    linkcolor=blue, citecolor=blue, urlcolor=blue,
    pdfauthor={Brieuc de La Fourniere}
}

% SPACING
\setstretch{1.2}
\setlength{\parskip}{0.4em}
\setlength{\parindent}{0pt}

% TITLE FORMATTING
\usepackage{titling}
\pretitle{\LARGE\bfseries}
\posttitle{\vspace{-0.4em}}
\preauthor{}
\postauthor{}
\predate{}
\postdate{}
\setlength{\droptitle}{-2.0em}

% CUSTOM COMMANDS — GIFT notation
\newcommand{\E}{\mathrm{E}}
\newcommand{\Gtwo}{\mathrm{G}_2}
\newcommand{\Kseven}{K_7}
\newcommand{\AdS}{\mathrm{AdS}}
\newcommand{\SU}{\mathrm{SU}}
\newcommand{\SO}{\mathrm{SO}}
\newcommand{\U}{\mathrm{U}}
\newcommand{\Sp}{\mathrm{Sp}}
\newcommand{\PSL}{\mathrm{PSL}}
\newcommand{\SM}{\mathrm{SM}}
\newcommand{\Tr}{\mathrm{Tr}}
\newcommand{\rank}{\mathrm{rank}}
\newcommand{\Vol}{\mathrm{Vol}}
\newcommand{\Hol}{\mathrm{Hol}}
\newcommand{\Ric}{\mathrm{Ric}}
\newcommand{\Rm}{\mathrm{Rm}}
\newcommand{\Scal}{\mathrm{Scal}}
\DeclareMathOperator{\diag}{diag}
\DeclareMathOperator{\cond}{cond}
\newcommand{\GIFT}{\textrm{GIFT}}
\newcommand{\CP}{\mathrm{CP}}
\newcommand{\CKM}{\mathrm{CKM}}
\newcommand{\PMNS}{\mathrm{PMNS}}
\newcommand{\EW}{\mathrm{EW}}
\newcommand{\Pl}{\mathrm{Pl}}
\newcommand{\DE}{\mathrm{DE}}
\newcommand{\DM}{\mathrm{DM}}
\newcommand{\KK}{\mathrm{KK}}
\newcommand{\NK}{\mathrm{NK}}
\newcommand{\TCS}{\mathrm{TCS}}
\newcommand{\PINN}{\mathrm{PINN}}
\newcommand{\btwo}{b_2}
\newcommand{\bthree}{b_3}
\newcommand{\typeI}{\textsc{I}}
\newcommand{\typeII}{\textsc{II}}
\newcommand{\typeIII}{\textsc{III}}
\newcommand{\typeIV}{\textsc{IV}}
\newcommand{\proven}{\textsc{Proven}}

% PDF bookmark safety
\pdfstringdefDisableCommands{%
  \def\textsubscript#1{#1}%
  \def\textsuperscript#1{#1}%
  \def\CP{CP}%
  \def\Gtwo{G2}%
  \def\Kseven{K7}%
  \def\GIFT{GIFT}%
  \def\E{E}%
  \def\SM{SM}%
}

% THEOREM ENVIRONMENTS
\newtheorem{theorem}{Theorem}[section]
\newtheorem{proposition}[theorem]{Proposition}
\newtheorem{lemma}[theorem]{Lemma}
\newtheorem{corollary}[theorem]{Corollary}
\theoremstyle{definition}
\newtheorem{definition}[theorem]{Definition}
\theoremstyle{remark}
\newtheorem{remark}[theorem]{Remark}
 after \documentclass

% ENCODING & FONTS
\usepackage[utf8]{inputenc}
\usepackage[T1]{fontenc}
\usepackage{lmodern}

% PAGE LAYOUT (Golden Ratio)
\usepackage[margin=1.618cm, top=2.618cm, bottom=2.618cm]{geometry}

% ESSENTIAL PACKAGES
\usepackage{float}
\usepackage{caption}
\usepackage{subcaption}
\usepackage{setspace}
\usepackage{fancyhdr}
\usepackage{xcolor}
\usepackage{hyperref}
\usepackage{csquotes}
\usepackage{amsmath}
\usepackage{amssymb}
\usepackage{amsthm}
\usepackage{mathtools}
\usepackage{booktabs}
\usepackage{longtable}
\usepackage{array}
\usepackage{multirow}
\usepackage{tikz}
\usepackage{graphicx}
\usepackage{listings}
\usepackage{enumitem}
% Unicode fallbacks (pandoc artifacts that survive markdown→LaTeX)
\DeclareUnicodeCharacter{00B0}{\ensuremath{^\circ}}
\DeclareUnicodeCharacter{221A}{\ensuremath{\sqrt{\,}}}
\DeclareUnicodeCharacter{2713}{\ensuremath{\checkmark}}
\DeclareUnicodeCharacter{2717}{\ensuremath{\times}}
\DeclareUnicodeCharacter{2070}{\ensuremath{^{0}}}
\DeclareUnicodeCharacter{00B9}{\ensuremath{^{1}}}
\DeclareUnicodeCharacter{00B2}{\ensuremath{^{2}}}
\DeclareUnicodeCharacter{00B3}{\ensuremath{^{3}}}
\DeclareUnicodeCharacter{2074}{\ensuremath{^{4}}}
\DeclareUnicodeCharacter{2075}{\ensuremath{^{5}}}
\DeclareUnicodeCharacter{2076}{\ensuremath{^{6}}}
\DeclareUnicodeCharacter{2077}{\ensuremath{^{7}}}
\DeclareUnicodeCharacter{2078}{\ensuremath{^{8}}}
\DeclareUnicodeCharacter{2079}{\ensuremath{^{9}}}
\DeclareUnicodeCharacter{2071}{\ensuremath{^{i}}}
\DeclareUnicodeCharacter{2080}{\ensuremath{_{0}}}
\DeclareUnicodeCharacter{2081}{\ensuremath{_{1}}}
\DeclareUnicodeCharacter{2082}{\ensuremath{_{2}}}
\DeclareUnicodeCharacter{2083}{\ensuremath{_{3}}}
\DeclareUnicodeCharacter{2084}{\ensuremath{_{4}}}
\DeclareUnicodeCharacter{2085}{\ensuremath{_{5}}}
\DeclareUnicodeCharacter{2086}{\ensuremath{_{6}}}
\DeclareUnicodeCharacter{2087}{\ensuremath{_{7}}}
\DeclareUnicodeCharacter{2088}{\ensuremath{_{8}}}
\DeclareUnicodeCharacter{2089}{\ensuremath{_{9}}}
\DeclareUnicodeCharacter{208A}{\ensuremath{_{+}}}
\DeclareUnicodeCharacter{208B}{\ensuremath{_{-}}}
\DeclareUnicodeCharacter{2192}{\ensuremath{\to}}
\DeclareUnicodeCharacter{2502}{|}
\DeclareUnicodeCharacter{25BC}{\ensuremath{\blacktriangledown}}
\DeclareUnicodeCharacter{2500}{\ensuremath{-}}
\DeclareUnicodeCharacter{2514}{\ensuremath{\llcorner}}
\DeclareUnicodeCharacter{251C}{\ensuremath{\vdash}}
\DeclareUnicodeCharacter{2026}{\ldots}
\DeclareUnicodeCharacter{2208}{\ensuremath{\in}}
\DeclareUnicodeCharacter{2229}{\ensuremath{\cap}}
\DeclareUnicodeCharacter{222A}{\ensuremath{\cup}}
\DeclareUnicodeCharacter{2295}{\ensuremath{\oplus}}
\DeclareUnicodeCharacter{2297}{\ensuremath{\otimes}}
\DeclareUnicodeCharacter{2245}{\ensuremath{\cong}}
\DeclareUnicodeCharacter{2248}{\ensuremath{\approx}}
\DeclareUnicodeCharacter{2260}{\ensuremath{\neq}}
\DeclareUnicodeCharacter{2261}{\ensuremath{\equiv}}
\DeclareUnicodeCharacter{2264}{\ensuremath{\leq}}
\DeclareUnicodeCharacter{2265}{\ensuremath{\geq}}
\DeclareUnicodeCharacter{00D7}{\ensuremath{\times}}
\DeclareUnicodeCharacter{00B1}{\ensuremath{\pm}}
\DeclareUnicodeCharacter{2225}{\ensuremath{\|}}
% Greek letters used in body text (pandoc passes them through unchanged)
\DeclareUnicodeCharacter{03B1}{\ensuremath{\alpha}}
\DeclareUnicodeCharacter{03B2}{\ensuremath{\beta}}
\DeclareUnicodeCharacter{03B3}{\ensuremath{\gamma}}
\DeclareUnicodeCharacter{03B4}{\ensuremath{\delta}}
\DeclareUnicodeCharacter{03B5}{\ensuremath{\varepsilon}}
\DeclareUnicodeCharacter{03B6}{\ensuremath{\zeta}}
\DeclareUnicodeCharacter{03B7}{\ensuremath{\eta}}
\DeclareUnicodeCharacter{03B8}{\ensuremath{\theta}}
\DeclareUnicodeCharacter{03B9}{\ensuremath{\iota}}
\DeclareUnicodeCharacter{03BA}{\ensuremath{\kappa}}
\DeclareUnicodeCharacter{03BB}{\ensuremath{\lambda}}
\DeclareUnicodeCharacter{03BC}{\ensuremath{\mu}}
\DeclareUnicodeCharacter{03BD}{\ensuremath{\nu}}
\DeclareUnicodeCharacter{03BE}{\ensuremath{\xi}}
\DeclareUnicodeCharacter{03C0}{\ensuremath{\pi}}
\DeclareUnicodeCharacter{03C1}{\ensuremath{\rho}}
\DeclareUnicodeCharacter{03C3}{\ensuremath{\sigma}}
\DeclareUnicodeCharacter{03C4}{\ensuremath{\tau}}
\DeclareUnicodeCharacter{03C5}{\ensuremath{\upsilon}}
\DeclareUnicodeCharacter{03C6}{\ensuremath{\varphi}}
\DeclareUnicodeCharacter{03C7}{\ensuremath{\chi}}
\DeclareUnicodeCharacter{03C8}{\ensuremath{\psi}}
\DeclareUnicodeCharacter{03C9}{\ensuremath{\omega}}
\DeclareUnicodeCharacter{0394}{\ensuremath{\Delta}}
\DeclareUnicodeCharacter{0398}{\ensuremath{\Theta}}
\DeclareUnicodeCharacter{039B}{\ensuremath{\Lambda}}
\DeclareUnicodeCharacter{03A0}{\ensuremath{\Pi}}
\DeclareUnicodeCharacter{03A3}{\ensuremath{\Sigma}}
\DeclareUnicodeCharacter{03A6}{\ensuremath{\Phi}}
\DeclareUnicodeCharacter{03A8}{\ensuremath{\Psi}}
\DeclareUnicodeCharacter{03A9}{\ensuremath{\Omega}}
\DeclareUnicodeCharacter{2200}{\ensuremath{\forall}}
\DeclareUnicodeCharacter{2203}{\ensuremath{\exists}}
\DeclareUnicodeCharacter{2207}{\ensuremath{\nabla}}
\DeclareUnicodeCharacter{2202}{\ensuremath{\partial}}
\DeclareUnicodeCharacter{2218}{\ensuremath{\circ}}
\DeclareUnicodeCharacter{2032}{\ensuremath{\prime}}
\DeclareUnicodeCharacter{2113}{\ensuremath{\ell}}
\DeclareUnicodeCharacter{210F}{\ensuremath{\hbar}}
\DeclareUnicodeCharacter{27E8}{\ensuremath{\langle}}
\DeclareUnicodeCharacter{27E9}{\ensuremath{\rangle}}
\DeclareUnicodeCharacter{2308}{\ensuremath{\lceil}}
\DeclareUnicodeCharacter{2309}{\ensuremath{\rceil}}
\DeclareUnicodeCharacter{230A}{\ensuremath{\lfloor}}
\DeclareUnicodeCharacter{230B}{\ensuremath{\rfloor}}
\DeclareUnicodeCharacter{222B}{\ensuremath{\int}}
\DeclareUnicodeCharacter{2211}{\ensuremath{\sum}}
\DeclareUnicodeCharacter{220F}{\ensuremath{\prod}}
\DeclareUnicodeCharacter{221E}{\ensuremath{\infty}}
\DeclareUnicodeCharacter{00B5}{\ensuremath{\mu}}
\DeclareUnicodeCharacter{2236}{\ensuremath{\colon}}
\DeclareUnicodeCharacter{2192}{\ensuremath{\to}}
\DeclareUnicodeCharacter{21A6}{\ensuremath{\mapsto}}
\DeclareUnicodeCharacter{21D2}{\ensuremath{\Rightarrow}}
\DeclareUnicodeCharacter{21D4}{\ensuremath{\Leftrightarrow}}
\DeclareUnicodeCharacter{2282}{\ensuremath{\subset}}
\DeclareUnicodeCharacter{2286}{\ensuremath{\subseteq}}
\DeclareUnicodeCharacter{2287}{\ensuremath{\supseteq}}
\DeclareUnicodeCharacter{222B}{\ensuremath{\int}}
\DeclareUnicodeCharacter{2205}{\ensuremath{\emptyset}}
\DeclareUnicodeCharacter{2299}{\ensuremath{\odot}}
\DeclareUnicodeCharacter{2220}{\ensuremath{\angle}}
\DeclareUnicodeCharacter{2135}{\ensuremath{\aleph}}
\DeclareUnicodeCharacter{207B}{\ensuremath{^{-}}}
\DeclareUnicodeCharacter{207A}{\ensuremath{^{+}}}
\DeclareUnicodeCharacter{220E}{\ensuremath{\blacksquare}}
\DeclareUnicodeCharacter{2061}{}
\DeclareUnicodeCharacter{2062}{}
\DeclareUnicodeCharacter{2063}{}
\DeclareUnicodeCharacter{2009}{\,}
\DeclareUnicodeCharacter{200A}{\,}
\DeclareUnicodeCharacter{200B}{}
\DeclareUnicodeCharacter{2212}{\ensuremath{-}}
\DeclareUnicodeCharacter{2194}{\ensuremath{\leftrightarrow}}
\DeclareUnicodeCharacter{2099}{\ensuremath{_{n}}}
\DeclareUnicodeCharacter{2098}{\ensuremath{_{m}}}
\DeclareUnicodeCharacter{2096}{\ensuremath{_{k}}}
\DeclareUnicodeCharacter{2095}{\ensuremath{_{h}}}
\DeclareUnicodeCharacter{2090}{\ensuremath{_{a}}}
\DeclareUnicodeCharacter{2091}{\ensuremath{_{e}}}
\DeclareUnicodeCharacter{2093}{\ensuremath{_{x}}}
\DeclareUnicodeCharacter{1D62}{\ensuremath{_{i}}}
\DeclareUnicodeCharacter{2C7C}{\ensuremath{_{j}}}
\DeclareUnicodeCharacter{2C7B}{\ensuremath{^{e}}}
\DeclareUnicodeCharacter{1D52}{\ensuremath{^{o}}}
\DeclareUnicodeCharacter{1D5}{\ensuremath{^{u}}}
\DeclareUnicodeCharacter{02E1}{\ensuremath{^{l}}}
\DeclareUnicodeCharacter{2295}{\ensuremath{\oplus}}
\DeclareUnicodeCharacter{27FA}{\ensuremath{\Longleftrightarrow}}
\DeclareUnicodeCharacter{27F9}{\ensuremath{\Longrightarrow}}
\DeclareUnicodeCharacter{2299}{\ensuremath{\odot}}
% Provide \tightlist used by pandoc-generated lists (no-op)
\providecommand{\tightlist}{\setlength{\itemsep}{0pt}\setlength{\parskip}{0pt}}
% pandoc longtable column specs: provide \real as passthrough and a
% `none` counter so that calc's \real{x} expansion path works.
\newcounter{none}
\providecommand{\real}[1]{#1}

% LISTINGS
\lstset{
    basicstyle=\small\ttfamily,
    breaklines=true, frame=single,
    keepspaces=true, showstringspaces=false,
    aboveskip=0.8em, belowskip=0.8em
}

% HEADER/FOOTER
\setlength{\headheight}{14pt}
\pagestyle{fancy}
\fancyhf{}
\fancyhead[R]{\thepage}
\renewcommand{\headrulewidth}{0.2pt}

% HYPERREF
\hypersetup{
    colorlinks=true,
    linkcolor=blue, citecolor=blue, urlcolor=blue,
    pdfauthor={Brieuc de La Fourniere}
}

% SPACING
\setstretch{1.2}
\setlength{\parskip}{0.4em}
\setlength{\parindent}{0pt}

% TITLE FORMATTING
\usepackage{titling}
\pretitle{\LARGE\bfseries}
\posttitle{\vspace{-0.4em}}
\preauthor{}
\postauthor{}
\predate{}
\postdate{}
\setlength{\droptitle}{-2.0em}

% CUSTOM COMMANDS — GIFT notation
\newcommand{\E}{\mathrm{E}}
\newcommand{\Gtwo}{\mathrm{G}_2}
\newcommand{\Kseven}{K_7}
\newcommand{\AdS}{\mathrm{AdS}}
\newcommand{\SU}{\mathrm{SU}}
\newcommand{\SO}{\mathrm{SO}}
\newcommand{\U}{\mathrm{U}}
\newcommand{\Sp}{\mathrm{Sp}}
\newcommand{\PSL}{\mathrm{PSL}}
\newcommand{\SM}{\mathrm{SM}}
\newcommand{\Tr}{\mathrm{Tr}}
\newcommand{\rank}{\mathrm{rank}}
\newcommand{\Vol}{\mathrm{Vol}}
\newcommand{\Hol}{\mathrm{Hol}}
\newcommand{\Ric}{\mathrm{Ric}}
\newcommand{\Rm}{\mathrm{Rm}}
\newcommand{\Scal}{\mathrm{Scal}}
\DeclareMathOperator{\diag}{diag}
\DeclareMathOperator{\cond}{cond}
\newcommand{\GIFT}{\textrm{GIFT}}
\newcommand{\CP}{\mathrm{CP}}
\newcommand{\CKM}{\mathrm{CKM}}
\newcommand{\PMNS}{\mathrm{PMNS}}
\newcommand{\EW}{\mathrm{EW}}
\newcommand{\Pl}{\mathrm{Pl}}
\newcommand{\DE}{\mathrm{DE}}
\newcommand{\DM}{\mathrm{DM}}
\newcommand{\KK}{\mathrm{KK}}
\newcommand{\NK}{\mathrm{NK}}
\newcommand{\TCS}{\mathrm{TCS}}
\newcommand{\PINN}{\mathrm{PINN}}
\newcommand{\btwo}{b_2}
\newcommand{\bthree}{b_3}
\newcommand{\typeI}{\textsc{I}}
\newcommand{\typeII}{\textsc{II}}
\newcommand{\typeIII}{\textsc{III}}
\newcommand{\typeIV}{\textsc{IV}}
\newcommand{\proven}{\textsc{Proven}}

% PDF bookmark safety
\pdfstringdefDisableCommands{%
  \def\textsubscript#1{#1}%
  \def\textsuperscript#1{#1}%
  \def\CP{CP}%
  \def\Gtwo{G2}%
  \def\Kseven{K7}%
  \def\GIFT{GIFT}%
  \def\E{E}%
  \def\SM{SM}%
}

% THEOREM ENVIRONMENTS
\newtheorem{theorem}{Theorem}[section]
\newtheorem{proposition}[theorem]{Proposition}
\newtheorem{lemma}[theorem]{Lemma}
\newtheorem{corollary}[theorem]{Corollary}
\theoremstyle{definition}
\newtheorem{definition}[theorem]{Definition}
\theoremstyle{remark}
\newtheorem{remark}[theorem]{Remark}

% pandoc >= 3.2.1 emits \pandocbounded around images; provide a fallback.
% Also clamp all images to \linewidth (natural-size figures overflowed the
% right margin -- the "truncated figures" flagged by the round-2 reviews).
% Kept here rather than in gift_preamble.tex, shared with the 3.4 build.
\providecommand{\pandocbounded}[1]{#1}
\makeatletter
\def\maxwidth{\ifdim\Gin@nat@width>\linewidth\linewidth\else\Gin@nat@width\fi}
\makeatother
\setkeys{Gin}{width=\maxwidth,keepaspectratio}
% md headings carry their own literal numbering ("3.2 What ..."); suppress
% the LaTeX section counters to avoid "0.2 1. Introduction" double numbering.
\setcounter{secnumdepth}{0}
% Loosen justification slightly: long technical identifiers (made breakable
% at build time via \allowbreak) still need extra stretch to avoid overfulls.
\emergencystretch=3em
\fancyhead[L]{The $K_7$ Framework v3.5: Supplement S3}
\hypersetup{pdftitle={The K7 Framework v3.5 Supplement S3: Observable Dataset}}

\graphicspath{{../../../figures/}}

\title{%
\LARGE\textbf{Supplement S3: Observable Dataset}\\[0.3em]
\large Complete Observable Dataset (95 entries) + Per-Sector Relation Tables%
}
\author{Brieuc de La Fourni\`ere \\ \textit{Independent researcher, Beaune, France}}
\date{July 2026 \quad\textbar\quad v3.5}

\begin{document}

\maketitle
\noindent\rule{\textwidth}{0.2pt}

\thispagestyle{fancy}

\section{S3.1 Introduction}\label{s3.1-introduction}

This supplement provides the complete observable dataset of the K₇ framework: 95 observables derived from a single G₂ metric (169 Chebyshev-Cholesky geometric parameters, zero freely adjustable physical parameters), organized by derivation type and cross-referenced with Lean formal verification. The dataset includes 2 BH remnant predictions from Pinčák et al.~2026 (main text ref. \cite{42}, classified Type IV) and 6 Type IV metric block eigenvalues (v3.4.29 ledger).

The 95 observables decompose as \textbf{33(I) + 19(II) + 21(III) + 22(IV)} across four types of increasing derivation complexity. All predictions trace to a single compact G₂ manifold K₇ with (b₂, b₃) = (21, 77) and E₈×E₈ gauge structure, 169 metric parameters, zero continuously adjustable physical parameters. The pair (b₂, b₃) = (21, 77) does not appear among the catalogued compact G₂ manifolds; the Joyce-Karigiannis Z₂³ route realizes it at the topological/lattice level (S1 §8.4), the companion scheme {[}E{]} discharges the datum-level analytic layer conditionally, and the global-analytic completion remains open (two-layer boundary: main §9). The certified metric provides computational evidence for a new G₂ manifold.

The dataset serves as the canonical reference for all framework predictions, superseding results previously scattered across internal working supplements. Every entry traces back to a specific computation, a JSON artifact, and (where applicable) a Lean theorem.

\begin{center}\rule{\linewidth}{0.5pt}\end{center}

\section{S3.2 Type Classification}\label{s3.2-type-classification}

Observables are classified into four types by derivation directness. The type assignment reflects the number of physical identification steps between the G₂ metric and the final prediction:

\begin{longtable}[]{@{}
  >{\raggedright\arraybackslash}p{(\columnwidth - 8\tabcolsep) * \real{0.1017}}
  >{\raggedright\arraybackslash}p{(\columnwidth - 8\tabcolsep) * \real{0.1186}}
  >{\raggedright\arraybackslash}p{(\columnwidth - 8\tabcolsep) * \real{0.2203}}
  >{\raggedright\arraybackslash}p{(\columnwidth - 8\tabcolsep) * \real{0.3051}}
  >{\raggedright\arraybackslash}p{(\columnwidth - 8\tabcolsep) * \real{0.2542}}@{}}
\toprule\noalign{}
\begin{minipage}[b]{\linewidth}\raggedright
Type
\end{minipage} & \begin{minipage}[b]{\linewidth}\raggedright
Count
\end{minipage} & \begin{minipage}[b]{\linewidth}\raggedright
Description
\end{minipage} & \begin{minipage}[b]{\linewidth}\raggedright
Derivation chain
\end{minipage} & \begin{minipage}[b]{\linewidth}\raggedright
Lean coverage
\end{minipage} \\
\midrule\noalign{}
\endhead
\bottomrule\noalign{}
\endlastfoot
I & 33 & Direct algebraic from structural constants & topology → formula → ratio & 33/33 (100\%) \\
II & 19 & One physical identification step & formula + VEV or scale → dimensional quantity & 0/19 (0\%) \\
III & 21 & Multi-step dynamical chains & metric → mechanism → observable & 14/21 (67\%) \\
IV & 22 & Structural/internal quantities & metric → diagnostic & 8/22 (36\%) \\
\end{longtable}

\textbf{Type I} (33 observables): Direct from metric geometry and topological invariants (b₂, b₃, dim(G₂), dim(E₈), etc.). These are dimensionless ratios expressed as algebraic combinations of integers. Examples: sin²θ\_\allowbreak{}W = 3/13, Q\_\allowbreak{}Koide = 2/3, m\_\allowbreak{}τ/m\_\allowbreak{}e = 3477. All 33 are Lean-certified. The two BH remnant topological predictions are classified as Type IV structural diagnostics (not part of the 33 Type I Lean-certified core).

\textbf{Type II} (19 observables): One physical identification step beyond Type I. Absolute masses from ratios × VEV (e.g., m\_\allowbreak{}u = (m\_\allowbreak{}u/m\_\allowbreak{}d) × m\_\allowbreak{}d\_\allowbreak{}exp), CKM magnitudes from Wolfenstein parametrization, extended quark ratios. From the \texttt{gift\_\allowbreak{}observables.\allowbreak{}csv} pipeline.

\textbf{Type III} (21 observables): Multi-step chains involving dynamical mechanisms: - wilson\_\allowbreak{}line non-adiabatic eigenvalue splitting (3 observables: raw lepton ratios) - S11 RGE gauge coupling running (4 observables: sin²θ\_\allowbreak{}W, α\_\allowbreak{}em⁻¹, α\_\allowbreak{}s at M\_\allowbreak{}Z, M\_\allowbreak{}GUT) - spectral 7D spectral geometry (5 observables: Weyl exponent, KK states, fiber channels) - gauge\_\allowbreak{}bundle gauge bundle data (4 observables: f\_\allowbreak{}IJ conditioning, α\_\allowbreak{}ratio, Yukawa rank) - instanton instanton volumes + combined wilson\_\allowbreak{}line+instanton (5 observables: ΔV's, combined ratios)

\textbf{Type IV} (22 observables): Structural and internal consistency quantities with no direct experimental comparison. These include topology diagnostics (b₂, b₃, χ(K₇)), Newton-Kantorovich certification values (h, distance, margin), Gram matrix conditioning (cond(G\_\allowbreak{}K3), cond(G\_\allowbreak{}77)), spectral counts, metric block eigenvalues (g\_\allowbreak{}ss = 19/6, g\_\allowbreak{}\{T²\} = 7/6, g\_\allowbreak{}K3 ≈ 64/77), and Pinčák et al.~2026 (main text ref. \cite{42}) topological diagnostics (N\_\allowbreak{}QNM = 98, M\_\allowbreak{}res, instanton suppression ×77). M\_\allowbreak{}res and N\_\allowbreak{}QNM are the two BH-remnant-related entries; they are classified here as structural diagnostics, not as Type I or Type III, because they require a physical scale identification (τ₀ = v\_\allowbreak{}EW) not derivable from topology alone. These quantities verify internal consistency rather than predict measurements directly.

\textbf{Derivation depth}: The type assignment correlates with derivation chain length. Type I observables require 1 algebraic step; Type II requires 2 steps (algebra + scale identification); Type III requires 3--5 steps (algebra + dynamics + calibration); Type IV involves diagnostic computation chains of variable length.

\begin{center}\rule{\linewidth}{0.5pt}\end{center}

\section{S3.3 Methodology}\label{s3.3-methodology}

\subsection{Experimental Sources}\label{experimental-sources}

All experimental reference values are drawn from four standard compilations:

\begin{longtable}[]{@{}
  >{\raggedright\arraybackslash}p{(\columnwidth - 4\tabcolsep) * \real{0.3333}}
  >{\raggedright\arraybackslash}p{(\columnwidth - 4\tabcolsep) * \real{0.4167}}
  >{\raggedright\arraybackslash}p{(\columnwidth - 4\tabcolsep) * \real{0.2500}}@{}}
\toprule\noalign{}
\begin{minipage}[b]{\linewidth}\raggedright
Source
\end{minipage} & \begin{minipage}[b]{\linewidth}\raggedright
Coverage
\end{minipage} & \begin{minipage}[b]{\linewidth}\raggedright
Date
\end{minipage} \\
\midrule\noalign{}
\endhead
\bottomrule\noalign{}
\endlastfoot
\textbf{PDG 2024} \cite{1} & Particle masses, gauge couplings, CKM elements & 2024 \\
\textbf{NuFIT 6.0} \cite{27} & Neutrino oscillation parameters (NO w/o SK) & Oct 2024 \\
\textbf{Planck PR4} (Tristram+ 2024) \cite{49} & Cosmological parameters: h, σ₈, Ω\_\allowbreak{}Λ & 2024 \\
\textbf{CODATA 2022} \cite{50} & Lepton mass ratios & 2022 \\
\end{longtable}

Selected NuFIT 6.0 values: sin²θ₂₃ = 0.561, δ\_\allowbreak{}CP = 177°. Selected Planck PR4 values: h = 0.6764, σ₈ = 0.807, Ω\_\allowbreak{}Λ = 0.6847. NuFIT 6.0 is the \emph{frozen dataset} for all deviation statistics in this supplement; the subsequent NuFIT 6.1 release (2025, δ\_\allowbreak{}CP = 207° +23/−20 NO w/o SK-atm) is tracked in the main paper §4.2 and moves the δ\_\allowbreak{}CP prediction inside the 1σ band.

\subsection{Structural Constants}\label{structural-constants}

The 20 structural constants from which all 95 observables derive:

\begin{longtable}[]{@{}
  >{\raggedright\arraybackslash}p{(\columnwidth - 6\tabcolsep) * \real{0.1071}}
  >{\raggedright\arraybackslash}p{(\columnwidth - 6\tabcolsep) * \real{0.3571}}
  >{\raggedright\arraybackslash}p{(\columnwidth - 6\tabcolsep) * \real{0.2500}}
  >{\raggedright\arraybackslash}p{(\columnwidth - 6\tabcolsep) * \real{0.2857}}@{}}
\toprule\noalign{}
\begin{minipage}[b]{\linewidth}\raggedright
\#
\end{minipage} & \begin{minipage}[b]{\linewidth}\raggedright
Constant
\end{minipage} & \begin{minipage}[b]{\linewidth}\raggedright
Value
\end{minipage} & \begin{minipage}[b]{\linewidth}\raggedright
Origin
\end{minipage} \\
\midrule\noalign{}
\endhead
\bottomrule\noalign{}
\endlastfoot
1 & dim(E₈) & 248 & E₈ Lie algebra \\
2 & rank(E₈) & 8 & Cartan subalgebra \\
3 & dim(G₂) & 14 & Aut(\ensuremath{\mathbb{O}}) \\
4 & dim(K₇) & 7 & Im(\ensuremath{\mathbb{O}}) \\
5 & b₂(K₇) & 21 & Input; confirmed by spectral analysis (Paper B {[}B{]}, §2.3). Any Mayer-Vietoris decomposition is conditional on building block identification (open problem). \\
6 & b₃(K₇) & 77 & Input; confirmed by spectral analysis (Paper B {[}B{]}, §2.3). Any Mayer-Vietoris decomposition is conditional on building block identification (open problem). \\
7 & N\_\allowbreak{}gen & 3 & Index theorem \\
8 & p₂ & 2 & dim(G₂)/dim(K₇) \\
9 & Weyl & 5 & Triple identity \\
10 & H* & 99 & b₂+b₃+1 \\
11 & dim(J₃(\ensuremath{\mathbb{O}})) & 27 & Jordan algebra \\
12 & dim(E₆) & 78 & E₆ Lie algebra \\
13 & dim(E₇) & 133 & E₇ Lie algebra \\
14 & dim(F₄) & 52 & Aut(J₃(\ensuremath{\mathbb{O}})) \\
15 & fund(E₇) & 56 & E₇ fundamental \\
16 & \textbar PSL(2,7)\textbar{} & 168 & Fano automorphisms \\
17 & D\_\allowbreak{}bulk & 11 & 4+7 \\
18 & α\_\allowbreak{}sum & 13 & rank(E₈)+Weyl \\
19 & det(g)\_\allowbreak{}den & 32 & b₂+dim(G₂)−N\_\allowbreak{}gen \\
20 & det(g)\_\allowbreak{}num & 65 & Weyl×α\_\allowbreak{}sum \\
\end{longtable}

\subsection{Consolidation Pipeline}\label{consolidation-pipeline}

The master dataset (\texttt{gift\_\allowbreak{}observables.\allowbreak{}json}) is generated by \texttt{observable\_\allowbreak{}dataset.\allowbreak{}py}, which collects the individual JSON results of the internal computation supplements and performs 5 verification checks:

\begin{enumerate}
\def\labelenumi{\arabic{enumi}.}
\tightlist
\item
  \textbf{Completeness}: All 95 entries present across 4 types
\item
  \textbf{Uniqueness}: No duplicate observable IDs
\item
  \textbf{Consistency}: Predicted values match source JSONs to machine precision
\item
  \textbf{Source tracing}: Every entry maps to a documented computation step
\item
  \textbf{Format validation}: JSON, CSV, and LaTeX outputs agree
\end{enumerate}

Each observable record contains: prediction (float64), fraction (exact symbolic), experimental value and uncertainty, experimental source, relative deviation (\%), Lean theorem name (if certified), and provenance (computation step).

\begin{center}\rule{\linewidth}{0.5pt}\end{center}

\section{S3.4 Complete Observable Table}\label{s3.4-complete-observable-table}

\textbf{Table S3.1}: Complete observable dataset (95 entries, v3.4.29 ledger). Columns: name, prediction, experimental value, relative deviation (\%), type (I/II/III/IV), source computation, Lean verification status.

{[}The full LaTeX table is provided in \texttt{data/\allowbreak{}gift\_\allowbreak{}observables\_\allowbreak{}latex.\allowbreak{}tex} and will be included at typesetting.{]}

\subsection{Per-Sector Relation Tables}\label{per-sector-relation-tables}

\emph{The per-sector Type I tables and their derivational prose moved here from main §4 at 3.5; the main paper keeps a representative summary (main §4.1) and the δ\_\allowbreak{}CP discussion (main §4.2). Conventions: PDG 2024 \cite{1}; NuFIT 6.0 frozen dataset \cite{27}; Planck PR4 \cite{49}.}

\subsubsection{Gauge sector}\label{gauge-sector}

\[\sin^2\theta_W = \frac{b_2}{b_3 + \dim(G_2)} = \frac{21}{91} = \frac{3}{13} = 0.230769\]

Experimental (PDG 2024) \cite{1}: 0.23122 ± 0.00004. Deviation: \textbf{0.195\%}. The numerator b₂ counts gauge moduli; the denominator b₃ + dim(G₂) counts matter plus holonomy degrees of freedom. The ratio measures gauge-matter coupling geometrically.

\[\alpha_s(M_Z) = \frac{\sqrt{2}}{\dim(G_2) - p_2} = \frac{\sqrt{2}}{12} = 0.11785\]

Experimental: 0.1180 ± 0.0009. Deviation: \textbf{0.126\%}.

\subsubsection{Lepton sector}\label{lepton-sector}

\textbf{Koide parameter.} The Koide relation among charged lepton masses, Q = (m\_\allowbreak{}e + m\_\allowbreak{}μ + m\_\allowbreak{}τ)/(√m\_\allowbreak{}e + √m\_\allowbreak{}μ + √m\_\allowbreak{}τ)² = 2/3, has resisted explanation since 1982. Using contemporary mass values it holds to six significant figures: Q\_\allowbreak{}exp = 0.666661 ± 0.000007. The K₇ framework provides Q\_\allowbreak{}Koide = dim(G₂)/b₂ = 14/21 = 2/3, an algebraic identity in two topological invariants; no fitting is involved.

\begin{longtable}[]{@{}
  >{\raggedright\arraybackslash}p{(\columnwidth - 4\tabcolsep) * \real{0.3846}}
  >{\raggedright\arraybackslash}p{(\columnwidth - 4\tabcolsep) * \real{0.3077}}
  >{\raggedright\arraybackslash}p{(\columnwidth - 4\tabcolsep) * \real{0.3077}}@{}}
\toprule\noalign{}
\begin{minipage}[b]{\linewidth}\raggedright
Approach
\end{minipage} & \begin{minipage}[b]{\linewidth}\raggedright
Result
\end{minipage} & \begin{minipage}[b]{\linewidth}\raggedright
Status
\end{minipage} \\
\midrule\noalign{}
\endhead
\bottomrule\noalign{}
\endlastfoot
Preon models (Koide 1982) & Q = 2/3 assumed & Circular \\
S₃ symmetry (various) & Q \textasciitilde{} 2/3 fitted & Approximate \\
\textbf{K₇ framework} & \textbf{Q = dim(G₂)/b₂ = 14/21 = 2/3} & \textbf{Algebraic identity} \\
\end{longtable}

Deviation: \textbf{0.0009\%}, the most precise agreement in the framework.

\textbf{Tau-electron mass ratio.} m\_\allowbreak{}τ/m\_\allowbreak{}e = dim(K₇) + 10 × dim(E₈) + 10 × H* = 7 + 2480 + 990 = 3477. Experimental: 3477.15 ± 0.05; deviation \textbf{0.004\%}. The integer 3477 = 3 × 19 × 61 = N\_\allowbreak{}gen × prime(8) × κ\_\allowbreak{}T⁻¹ factorizes into framework constants.

\textbf{Muon-electron mass ratio.} m\_\allowbreak{}μ/m\_\allowbreak{}e = dim(J₃(\ensuremath{\mathbb{O}}))\^{}φ = 27\^{}φ = 207.01, with φ = (1+√5)/2 from the McKay correspondence E₈ ↔ 2I (binary icosahedral group). Experimental: 206.768; deviation \textbf{0.118\%}. See the golden-ratio caveat in main §4.1: φ is the only non-integer input among the 33 Type I relations, and its status is weaker than the integer-ratio derivations (S2 §11).

\subsubsection{Quark sector}\label{quark-sector}

\[\frac{m_s}{m_d} = p_2^2 \times \text{Weyl} = 4 \times 5 = 20\]

Experimental (PDG 2024): 20.0 ± 1.0. Deviation: \textbf{0.00\%}.

\[\frac{m_b}{m_t} = \frac{b_0}{2b_2} = \frac{1}{42}\]

Experimental: 0.024 ± 0.001. Deviation: \textbf{0.79\%}. The constant 42 = p₂ × N\_\allowbreak{}gen × dim(K₇) = 2 × 3 × 7 is a structural invariant (not to be confused with χ(K₇) = 0, which vanishes for any compact odd-dimensional manifold).

\subsubsection{Neutrino mixing angles}\label{neutrino-mixing-angles}

\begin{longtable}[]{@{}lllll@{}}
\toprule\noalign{}
Angle & Formula & Prediction & NuFIT 6.0 \cite{27} & Dev. \\
\midrule\noalign{}
\endhead
\bottomrule\noalign{}
\endlastfoot
θ₁₂ & arctan(√(δ/γ\_\allowbreak{}GIFT)) & 33.40° & 33.41 ± 0.75° & 0.03\% \\
θ₁₃ & π/b₂ & 8.57° & 8.54 ± 0.12° & 0.37\% \\
θ₂₃ & arcsin((b₃ − p₂)/H*) & 49.25° & 49.3 ± 1.0° & 0.10\% \\
\end{longtable}

Auxiliary parameters: δ = 2π/Weyl² = 2π/25 and γ\_\allowbreak{}GIFT (historical symbol name) = (2 × rank(E₈) + 5 × H*)/(10 × dim(G₂) + 3 × dim(E₈)) = 511/884. The δ\_\allowbreak{}CP prediction is discussed in main §4.2.

\subsubsection{Higgs sector}\label{higgs-sector}

\[\lambda_H = \frac{\sqrt{\dim(G_2) + N_{gen}}}{\det(g)_{den}} = \frac{\sqrt{17}}{32} = 0.1289\]

The Higgs self-coupling combines holonomy dimension with generation count, normalized by the metric determinant scale. Experimental (PDG 2024): 0.129 ± 0.001. Deviation: \textbf{0.12\%}. A TCS alternative λ\_\allowbreak{}H = b₂(M₁)/(b₃+b₂(M₂)) = 11/87 = 0.1264 is purely rational (see S2 §17).

\subsubsection{Boson mass ratios}\label{boson-mass-ratios}

\begin{longtable}[]{@{}
  >{\raggedright\arraybackslash}p{(\columnwidth - 8\tabcolsep) * \real{0.2553}}
  >{\raggedright\arraybackslash}p{(\columnwidth - 8\tabcolsep) * \real{0.1915}}
  >{\raggedright\arraybackslash}p{(\columnwidth - 8\tabcolsep) * \real{0.1277}}
  >{\raggedright\arraybackslash}p{(\columnwidth - 8\tabcolsep) * \real{0.2979}}
  >{\raggedright\arraybackslash}p{(\columnwidth - 8\tabcolsep) * \real{0.1277}}@{}}
\toprule\noalign{}
\begin{minipage}[b]{\linewidth}\raggedright
Observable
\end{minipage} & \begin{minipage}[b]{\linewidth}\raggedright
Formula
\end{minipage} & \begin{minipage}[b]{\linewidth}\raggedright
Prediction
\end{minipage} & \begin{minipage}[b]{\linewidth}\raggedright
Experimental
\end{minipage} & \begin{minipage}[b]{\linewidth}\raggedright
Dev.
\end{minipage} \\
\midrule\noalign{}
\endhead
\bottomrule\noalign{}
\endlastfoot
m\_\allowbreak{}H/m\_\allowbreak{}W & (N\_\allowbreak{}gen + dim(E₆))/dim(F₄) = 81/52 & 1.5577 & 1.558 ± 0.002 & 0.02\% \\
m\_\allowbreak{}W/m\_\allowbreak{}Z & (2b₂ - Weyl)/(2b₂) = 37/42 & 0.8810 & 0.8815 ± 0.0002 & 0.06\% \\
m\_\allowbreak{}H/m\_\allowbreak{}t & fund(E₇)/b₃ = 56/77 & 0.7273 & 0.725 ± 0.003 & 0.31\% \\
\end{longtable}

\subsubsection{CKM matrix}\label{ckm-matrix}

\begin{longtable}[]{@{}
  >{\raggedright\arraybackslash}p{(\columnwidth - 8\tabcolsep) * \real{0.2553}}
  >{\raggedright\arraybackslash}p{(\columnwidth - 8\tabcolsep) * \real{0.1915}}
  >{\raggedright\arraybackslash}p{(\columnwidth - 8\tabcolsep) * \real{0.1277}}
  >{\raggedright\arraybackslash}p{(\columnwidth - 8\tabcolsep) * \real{0.2979}}
  >{\raggedright\arraybackslash}p{(\columnwidth - 8\tabcolsep) * \real{0.1277}}@{}}
\toprule\noalign{}
\begin{minipage}[b]{\linewidth}\raggedright
Observable
\end{minipage} & \begin{minipage}[b]{\linewidth}\raggedright
Formula
\end{minipage} & \begin{minipage}[b]{\linewidth}\raggedright
Prediction
\end{minipage} & \begin{minipage}[b]{\linewidth}\raggedright
Experimental
\end{minipage} & \begin{minipage}[b]{\linewidth}\raggedright
Dev.
\end{minipage} \\
\midrule\noalign{}
\endhead
\bottomrule\noalign{}
\endlastfoot
sin²(θ₁₂ CKM) & fund(E₇)/dim(E₈) = 56/248 & 0.2258 & 0.2250 ± 0.0006 & 0.36\% \\
A\_\allowbreak{}Wolfenstein & (Weyl + dim(E₆))/H* = 83/99 & 0.838 & 0.836 ± 0.015 & 0.29\% \\
sin²(θ₂₃ CKM) & dim(K₇)/PSL(2,7) = 7/168 & 0.0417 & 0.0412 ± 0.0008 & 1.13\% \\
\end{longtable}

The Cabibbo angle emerges from the ratio of the E₇ fundamental representation to E₈ dimension.

\subsubsection{Cosmological observables}\label{cosmological-observables}

\begin{longtable}[]{@{}
  >{\raggedright\arraybackslash}p{(\columnwidth - 8\tabcolsep) * \real{0.2553}}
  >{\raggedright\arraybackslash}p{(\columnwidth - 8\tabcolsep) * \real{0.1915}}
  >{\raggedright\arraybackslash}p{(\columnwidth - 8\tabcolsep) * \real{0.1277}}
  >{\raggedright\arraybackslash}p{(\columnwidth - 8\tabcolsep) * \real{0.2979}}
  >{\raggedright\arraybackslash}p{(\columnwidth - 8\tabcolsep) * \real{0.1277}}@{}}
\toprule\noalign{}
\begin{minipage}[b]{\linewidth}\raggedright
Observable
\end{minipage} & \begin{minipage}[b]{\linewidth}\raggedright
Formula
\end{minipage} & \begin{minipage}[b]{\linewidth}\raggedright
Prediction
\end{minipage} & \begin{minipage}[b]{\linewidth}\raggedright
Experimental
\end{minipage} & \begin{minipage}[b]{\linewidth}\raggedright
Dev.
\end{minipage} \\
\midrule\noalign{}
\endhead
\bottomrule\noalign{}
\endlastfoot
Ω\_\allowbreak{}DM/Ω\_\allowbreak{}b & (1 + 2b₂)/rank(E₈) = 43/8 & 5.375 & 5.375 ± 0.1 & 0.00\% \\
n\_\allowbreak{}s & ζ(11)/ζ(5) & 0.9649 & 0.9649 ± 0.0042 & 0.004\% \\
h (Hubble) & (PSL(2,7) - 1)/dim(E₈) = 167/248 & 0.6734 & 0.674 ± 0.005 & 0.09\% \\
Ω\_\allowbreak{}b/Ω\_\allowbreak{}m & Weyl/det(g)\_\allowbreak{}den = 5/32 & 0.1562 & 0.157 ± 0.003 & 0.16\% \\
σ₈ & (p₂ + 32)/(2b₂) = 34/42 & 0.8095 & 0.811 ± 0.006 & 0.18\% \\
Ω\_\allowbreak{}DE & ln(2)×(b₂+b₃)/H* = ln(2)×98/99 & 0.6861 & 0.6847 ± 0.005 & 0.21\% \\
Y\_\allowbreak{}p & (1 + dim(G₂))/κ\_\allowbreak{}T⁻¹ = 15/61 & 0.2459 & 0.245 ± 0.003 & 0.37\% \\
\end{longtable}

The dark-to-baryonic matter ratio Ω\_\allowbreak{}DM/Ω\_\allowbreak{}b = 43/8 is exact. The structural invariant 2b₂ = 42 that gives m\_\allowbreak{}b/m\_\allowbreak{}t = 1/42 also determines this cosmological ratio, connecting quark physics to large-scale structure through K₇ geometry.

\subsection{Sector-by-Sector Performance}\label{sector-by-sector-performance}

The 66 observables with experimental comparison (Types I+II+III) organize into sectors of varying precision:

\textbf{Best-performing sectors} (mean dev \textless{} 0.5\%): - \textbf{Lepton} (3 obs): mean 0.04\%: Q\_\allowbreak{}Koide (0.001\%), m\_\allowbreak{}τ/m\_\allowbreak{}e (0.004\%), m\_\allowbreak{}μ/m\_\allowbreak{}e (0.12\%) - \textbf{Electroweak} (3 obs): mean 0.11\%, α⁻¹ (0.002\%), α\_\allowbreak{}s (0.126\%), sin²θ\_\allowbreak{}W (0.20\%) - \textbf{Boson} (3 obs): mean 0.13\%, m\_\allowbreak{}H/m\_\allowbreak{}W (0.02\%), m\_\allowbreak{}W/m\_\allowbreak{}Z (0.06\%), m\_\allowbreak{}H/m\_\allowbreak{}t (0.31\%) - \textbf{Cosmology} (7 obs): mean 0.15\%, n\_\allowbreak{}s (0.004\%), h (0.09\%), Ω\_\allowbreak{}b/Ω\_\allowbreak{}m (0.16\%), σ₈ (0.18\%), Ω\_\allowbreak{}DE (0.21\%), Y\_\allowbreak{}p (0.37\%), Ω\_\allowbreak{}DM/Ω\_\allowbreak{}b (0.00\%) - \textbf{Quark} (4 obs): mean 0.24\%, m\_\allowbreak{}s/m\_\allowbreak{}d (0.00\%), m\_\allowbreak{}u/m\_\allowbreak{}d (0.05\%), m\_\allowbreak{}c/m\_\allowbreak{}s (0.12\%), m\_\allowbreak{}b/m\_\allowbreak{}t (0.79\%) - \textbf{PMNS} (4 obs): mean 0.29\%, θ₁₂ (0.03\%), θ₂₃ (0.10\%), sin²θ₁₂ (0.23\%), sin²θ₁₃ (0.81\%)

\textbf{Good sectors} (mean dev 0.5--3\%): - \textbf{CKM} (3 obs): mean 0.59\%: A\_\allowbreak{}Wolf (0.29\%), sin²θ₁₂ (0.36\%), sin²θ₂₃ (1.13\%) - \textbf{Gauge running} (4 obs): mean 2.3\%, sin²θ\_\allowbreak{}W(RGE) (2.78\%), α\_\allowbreak{}em⁻¹(RGE) (2.53\%), α\_\allowbreak{}s(RGE) (3.7\%), M\_\allowbreak{}GUT (0.00\%)

\textbf{Challenging sectors} (mean dev \textgreater{} 3\%): - \textbf{Instanton} (2 obs): mean 10.9\%: ΔV(e-τ) (5.9\%), ΔV(e-μ) (15.9\%); both reduced to sub-percent by the combined pipeline (§S3.5) - \textbf{Bundle} (1 obs): κ(gauge) (4.7\%), classified with the gauge-bundle diagnostics as in main §5.5

\textbf{B-test consistency} (derived, not counted as independent observable): The MSSM B parameter B = (α₁⁻¹−α₂⁻¹)/(α₂⁻¹−α₃⁻¹) evaluates to 1.4033 using the framework couplings (0.23\% from the theoretical 7/5), closer than the purely experimental B = 1.3948 (0.37\%). The identity B = 7/5 is equivalent to α\_\allowbreak{}em⁻¹(M\_\allowbreak{}Z) = 91√2 = dim(Λ²\ensuremath{\mathfrak{g}}₂)·√2 and implies the holonomy sequence α₁⁻¹:α₂⁻¹:α₃⁻¹ = dim(G₂):dim(K₇):p₂ = 14:7:2. See §5.4 of the main text.

\subsection{Exact Matches}\label{exact-matches}

11 observables achieve deviations below 0.01\% (effectively exact):

\begin{longtable}[]{@{}llll@{}}
\toprule\noalign{}
Observable & Prediction & Exp. & Dev. \\
\midrule\noalign{}
\endhead
\bottomrule\noalign{}
\endlastfoot
α⁻¹ & 137.033 & 137.036 & 0.002\% \\
Q\_\allowbreak{}Koide & 2/3 & 0.666661 & 0.001\% \\
m\_\allowbreak{}τ/m\_\allowbreak{}e & 3477 & 3477.15 & 0.004\% \\
m\_\allowbreak{}s/m\_\allowbreak{}d & 20 & 20.0 & 0.000\% \\
n\_\allowbreak{}s & 0.9649 & 0.9649 & 0.004\% \\
m\_\allowbreak{}u & 2.16 & 2.16 & 0.00\% \\
\textbar V\_\allowbreak{}tb\textbar{} & 0.999 & 0.999 & 0.00\% \\
m\_\allowbreak{}c & 1.27 & 1.27 & 0.00\% \\
M\_\allowbreak{}GUT & 2×10¹⁶ & 2×10¹⁶ & 0.00\% \\
α\_\allowbreak{}GUT⁻¹ & 25.3 & 25.3 & 0.00\% \\
rank(Y) & 3 & 3 & 0.00\% \\
\end{longtable}

Note: δ\_\allowbreak{}CP canonical prediction is 197° (11.3\% from the frozen NuFIT 6.0 central; 4.8\% from the NuFIT 6.1 central, inside 1σ); a compactification factor 62/69 brings it to 177.014° vs the 6.0 central (0.008\%) but is not adopted, see main §4.2. m\_\allowbreak{}c/m\_\allowbreak{}s deviation is 0.12\% and its Type I product with m\_\allowbreak{}s/m\_\allowbreak{}d gives m\_\allowbreak{}c/m\_\allowbreak{}d = 234.3 vs the PDG scale-common ratio 234 (0.12\%), so neither is listed among the exact matches above; the mixed-scale value 272 that would arise from combining m\_\allowbreak{}c(m\_\allowbreak{}c) and m\_\allowbreak{}d(2 GeV) violates the §4 convention (all quark ratios at the PDG reference scales) and is not tracked.

\subsection{Known Outliers}\label{known-outliers}

Two observables deviate by more than 5\%, each with an identified physical origin:

\textbf{δ\_\allowbreak{}CP}: The framework prediction is 197° (pure topological: 7 × 14 + 99), deviating 11.3\% from the frozen NuFIT 6.0 central (177° ± 20°, edge of 1σ) and 4.8\% from the NuFIT 6.1 central (207° +23/−20, inside 1σ). PSLQ analysis identifies a compactification factor 62/69 (§4.2 of main paper), documented as a structural observation but not adopted: the canonical prediction remains 197°. The experimental central value drifted from 197° (NuFIT 5.2) to 177° (NuFIT 6.0) to 207° (NuFIT 6.1), bracketing the prediction. Falsifiable by DUNE (first beam targeted 2031; physics run late 2030s--2040s).

\textbf{α\_\allowbreak{}s(RGE) = +3.7\%}: G₂-MSSM split-spectrum matching (MSSM above M\_\allowbreak{}squark = 3165 GeV, SM + gauginos below (b₃ = −5)) gives α\_\allowbreak{}s = 0.1224 (exp 0.1180, dev 3.7\%). A naive degenerate-spectrum treatment would give 12.1\%. The topological value α\_\allowbreak{}s = √2/12 = 0.11785 (Type I) matches experiment to 0.126\%.

\textbf{ΔV(e-μ) = +15.9\%}: The instanton volume assignment ΔV(e-μ) = 3.271 versus target ln(16.82) = 2.823 reflects the optimal assignment problem for 57 associative cycles. The combined wilson\_\allowbreak{}line+instanton pipeline with α = e\^{}K (§S3.5) reduces this to 0.75\%.

\textbf{κ(gauge) = +4.7\%}: The gauge kinetic function conditioning cond(f\_\allowbreak{}IJ) = 1.047 measures departure from exact universality. The deviation from 1.000 reflects K₇ geometry at the percent level.

\subsection{Summary Statistics}\label{summary-statistics}

\begin{longtable}[]{@{}ll@{}}
\toprule\noalign{}
Metric & Value \\
\midrule\noalign{}
\endhead
\bottomrule\noalign{}
\endlastfoot
Total observables & 95 \\
Type breakdown & 33(I) + 19(II) + 21(III) + 22(IV) \\
With experimental comparison & 66 \\
Exact matches (dev \textless{} 0.01\%) & 11 \\
Within 1\% of experiment & 53 \\
Within 5\% of experiment & 63 \\
Lean-certified & 55 \\
Structural constants & 20 \\
Metric parameters & 169 \\
\end{longtable}

\begin{center}\rule{\linewidth}{0.5pt}\end{center}

\section{S3.5 The Combined wilson\_\allowbreak{}line+instanton Pipeline}\label{s3.5-the-combined-wilson_lineinstanton-pipeline}

The lepton mass hierarchy (combined lepton ratio observables) uses two complementary geometric mechanisms. The complete derivation is in §6 of the main text; results are summarised here for cross-reference.

Raw wilson\_\allowbreak{}line (K3 fiber eigenvalue splitting at c = 0.452): m\_\allowbreak{}τ/m\_\allowbreak{}e = 3403 (2.1\%), m\_\allowbreak{}τ/m\_\allowbreak{}μ = 16.54 (1.7\%), m\_\allowbreak{}μ/m\_\allowbreak{}e = 205.7 (0.5\%). Raw instanton (associative 3-cycle volumes, 57 cycles): ΔV(e-τ) = 8.633 (5.9\%), ΔV(e-μ) = 3.271 (15.9\%). Combined with α = e\^{}\{K₀\} = V̂\^{}\{−3\} = 0.002763 (Kähler potential of K₇, zero free parameters):

\begin{itemize}
\tightlist
\item
  m\_\allowbreak{}τ/m\_\allowbreak{}e = 3485 (exp 3477, \textbf{0.23\%})
\item
  m\_\allowbreak{}τ/m\_\allowbreak{}μ = 16.69 (exp 16.82, \textbf{0.75\%})
\item
  m\_\allowbreak{}μ/m\_\allowbreak{}e = 208.8 (exp 206.8, \textbf{0.98\%})
\end{itemize}

Certified in Lean: \texttt{Associative\allowbreak{}Volumes.\allowbreak{}lean} (14 conjuncts, 0 sorry). Full mechanism description: main §6.4.

\begin{center}\rule{\linewidth}{0.5pt}\end{center}

\section{S3.6 Sensitivity Cross-Reference}\label{s3.6-sensitivity-cross-reference}

The complete sensitivity analysis is in §7 of the main text. Key results for cross-reference:

\begin{itemize}
\tightlist
\item
  \textbf{Effective DOF}: r\_\allowbreak{}eff = 15.53 (SVD of 20×33 constant-usage Jacobian)
\item
  \textbf{Overdetermination}: 33/15.53 = 2.13× (more constraints than free dimensions)
\item
  \textbf{Cross-coupling}: 155/528 observable pairs share a structural constant; all form one connected component
\item
  \textbf{P(coincidence, uniform)} = 10⁻³⁴⁶ (χ² + Fisher combined; archived diagnostic, see Supplement S4)
\item
  \textbf{P(coincidence, algebraic)} = 10⁻¹³³ (4.2M random formulas from same 20 constants; archived diagnostic, see Supplement S4). These joint-probability figures are internal-consistency diagnostics, superseded as headline methodology by the Sieve reading (main §7).
\end{itemize}

Figures: \texttt{fig\_\allowbreak{}sensitivity\_\allowbreak{}heatmap.\allowbreak{}png}, \texttt{fig\_\allowbreak{}constant\_\allowbreak{}usage.\allowbreak{}png}, \texttt{fig\_\allowbreak{}observable\_\allowbreak{}correlations.\allowbreak{}png}, \texttt{fig\_\allowbreak{}mc\_\allowbreak{}per\_\allowbreak{}observable.\allowbreak{}png}.

\begin{center}\rule{\linewidth}{0.5pt}\end{center}

\section{S3.7 Lean Cross-References}\label{s3.7-lean-cross-references}

Of the 95 observables, \textbf{55 are formally verified} in Lean 4 across the following certificate files:

\begin{longtable}[]{@{}
  >{\raggedright\arraybackslash}p{(\columnwidth - 8\tabcolsep) * \real{0.2200}}
  >{\raggedright\arraybackslash}p{(\columnwidth - 8\tabcolsep) * \real{0.1600}}
  >{\raggedright\arraybackslash}p{(\columnwidth - 8\tabcolsep) * \real{0.2000}}
  >{\raggedright\arraybackslash}p{(\columnwidth - 8\tabcolsep) * \real{0.2200}}
  >{\raggedright\arraybackslash}p{(\columnwidth - 8\tabcolsep) * \real{0.2000}}@{}}
\toprule\noalign{}
\begin{minipage}[b]{\linewidth}\raggedright
Lean File
\end{minipage} & \begin{minipage}[b]{\linewidth}\raggedright
Axioms
\end{minipage} & \begin{minipage}[b]{\linewidth}\raggedright
Theorems
\end{minipage} & \begin{minipage}[b]{\linewidth}\raggedright
Conjuncts
\end{minipage} & \begin{minipage}[b]{\linewidth}\raggedright
Coverage
\end{minipage} \\
\midrule\noalign{}
\endhead
\bottomrule\noalign{}
\endlastfoot
\texttt{Foundations.\allowbreak{}lean} & 0* & n/a & 39 & Type I: metric, torsion, topology \\
\texttt{Predictions.\allowbreak{}lean} & 0* & n/a & 56 & Types I--III: couplings, masses, mixing \\
\texttt{Spectral.\allowbreak{}lean} & 0* & n/a & 45 & Types III-IV: KK spectrum, Weyl law \\
\texttt{Metric\allowbreak{}Eigenvalues.\allowbreak{}lean} & 0 & n/a & 15 & Metric fractions: g\_\allowbreak{}ss=19/6, g\_\allowbreak{}T²=7/6, γ²=24π²/7 (T² Laplacian) \\
\texttt{Spectral\allowbreak{}Invariants.\allowbreak{}lean} & 0 & n/a & 10 & Heat kernel, ζ'(0), b₁=0 \\
\texttt{Compactification\allowbreak{}Correction.\allowbreak{}lean} & 0 & n/a & 6 & δ\_\allowbreak{}CP compactification factor \\
\texttt{TCSGauge\allowbreak{}Breaking.\allowbreak{}lean} & 0 & 14 & 10 & Type III: gauge breaking chain \\
\texttt{Gauge\allowbreak{}Bundle\allowbreak{}Data.\allowbreak{}lean} & 0 & 12 & 11 & Type III: bundle universality \\
\texttt{Associative\allowbreak{}Volumes.\allowbreak{}lean} & 0 & 19 & 14 & Type III: instanton hierarchy \\
\texttt{Computed\allowbreak{}Weyl\allowbreak{}Law.\allowbreak{}lean} & 0 & 8 & 7 & Type III: 7D spectral geometry \\
\end{longtable}

*4 external data-package axioms (of 15 classified in the A--F taxonomy of the v3.4.29 release audit; see main §8.1 for the reconciliation). All substantive: standard theorems (Cheeger inequality) + geometric structure (TCS spectral bounds) + physical inputs (literature package: CGN+Joyce). G₂ group structure proven by \texttt{native\_\allowbreak{}decide}.

\textbf{Total certificate}: 213 conjuncts (top-level conjunctions of the master certificates: 39 + 56 + 45 + 15 + 10 + 6 + 10 + 11 + 14 + 7), 0 \texttt{sorry}, 146 .lean files under \texttt{GIFT/\allowbreak{}} (140 core + 6 generated), 8394 build jobs (Lean 4.29.0).

\subsection{Coverage by Type}\label{coverage-by-type}

\begin{longtable}[]{@{}llll@{}}
\toprule\noalign{}
Type & Certified & Total & Coverage \\
\midrule\noalign{}
\endhead
\bottomrule\noalign{}
\endlastfoot
\textbf{I} & 33 & 33 & 100\% \\
\textbf{II} & 0 & 19 & 0\% \\
\textbf{III} & 14 & 21 & 67\% \\
\textbf{IV} & 8 & 22 & 36\% \\
\textbf{Total} & \textbf{55} & \textbf{95} & \textbf{58\%} \\
\end{longtable}

\textbf{Type II uncertified}: The 19 Type II observables derive from the CSV pipeline (absolute masses from ratio × VEV, CKM magnitudes from Wolfenstein). These involve dimensional quantities and intermediate data that have not yet been axiomatized in Lean. The underlying Type I ratios from which they derive are all certified.

\textbf{Type IV partial}: 8 of 22 structural quantities are certified (Betti numbers, holonomy dimension, NK certification bounds). The remaining 14 involve spectral counts, conditioning numbers, metric eigenvalues, and BH remnant estimates that require floating-point Lean infrastructure not yet developed.

\begin{center}\rule{\linewidth}{0.5pt}\end{center}

\section{S3.8 Data Availability}\label{s3.8-data-availability}

The complete dataset is generated deterministically by the \texttt{observable\_\allowbreak{}dataset.\allowbreak{}py} + \texttt{deviation\_\allowbreak{}systematics.\allowbreak{}py} pipeline from the frozen inputs and is available in three formats:

\begin{itemize}
\tightlist
\item
  \textbf{JSON}: \texttt{gift\_\allowbreak{}observables.\allowbreak{}json} (machine-readable, full metadata per observable)
\item
  \textbf{CSV}: \texttt{gift\_\allowbreak{}observables.\allowbreak{}csv} (tabular, 95 rows × 13 columns, for analysis)
\item
  \textbf{LaTeX}: \texttt{gift\_\allowbreak{}observables\_\allowbreak{}latex.\allowbreak{}tex} (publication-ready table)
\end{itemize}

The dataset, the generator scripts, and the mechanism-selection protocols referenced in main §6 (search spaces, objective function, cycle-assignment selection) are available from the author on request pending public deposit. The spectral companion artifacts of {[}B{]} are archived separately on Zenodo (DOI: \href{https://doi.org/10.5281/zenodo.19893371}{10.5281/zenodo.19893371}).

\begin{center}\rule{\linewidth}{0.5pt}\end{center}


\bigskip
% ============================================
% GIFT v3.5 — Shared Bibliography
% ============================================
% Usage: % ============================================
% GIFT v3.5 — Shared Bibliography
% ============================================
% Usage: % ============================================
% GIFT v3.5 — Shared Bibliography
% ============================================
% Usage: \input{gift_3.5_bib} where the references section should appear.
% Sequential numbering [1]..[50]. Round-1 + round-2 IA-review updates applied
% (2026-07-07): [13] Kovalev 2005, [26] T2K/NOvA Nature 646,
% [27] NuFit 6.0 + 6.1, [45] Garbagnati-Sarti 2007, [46] LEPSUSYWG,
% [47] Fermilab DUNE 2031, [48] Chen-Hong (re-added, cited in S1), [49] Planck PR4, [50] CODATA.

\begin{thebibliography}{99}

\bibitem{1} Particle Data Group, ``Review of Particle Physics,'' \textit{Phys.\ Rev.\ D} \textbf{110}, 030001 (2024).

\bibitem{2} S.~Weinberg, ``Implications of dynamical symmetry breaking,'' \textit{Phys.\ Rev.\ D} \textbf{13}, 974 (1976).

\bibitem{3} Planck Collaboration, ``Planck 2018 results.\ VI.\ Cosmological parameters,'' \textit{Astron.\ Astrophys.} \textbf{641}, A6 (2020).

\bibitem{4} A.~G.~Riess et al., ``A comprehensive measurement of the local value of the Hubble constant,'' \textit{ApJL} \textbf{934}, L7 (2022).

\bibitem{5} C.~D.~Froggatt, H.~B.~Nielsen, ``Hierarchy of quark masses, Cabibbo angles and CP violation,'' \textit{Nucl.\ Phys.\ B} \textbf{147}, 277 (1979).

\bibitem{6} Y.~Koide, ``Fermion-boson two-body model of quarks and leptons and Cabibbo mixing,'' \textit{Lett.\ Nuovo Cim.} \textbf{34}, 201 (1982).

\bibitem{7} C.~Furey, \textit{Standard Model Physics from an Algebra?}, PhD thesis, University of Waterloo (2015).

\bibitem{8} N.~Furey, M.~J.~Hughes, ``One generation of standard model Weyl representations as a single copy of $\mathbb{R} \otimes \mathbb{C} \otimes \mathbb{H} \otimes \mathbb{O}$,'' \textit{Phys.\ Lett.\ B} \textbf{831}, 137186 (2022).

\bibitem{9} R.~Wilson, ``$E_8$ and Standard Model plus gravity,'' arXiv:2404.18938 [hep-th] (2024).

\bibitem{10} T.~P.~Singh et al., ``An $E_8 \otimes E_8$ unification of the Standard Model with pre-gravitation,'' arXiv:2206.06911v3 (2024).

\bibitem{11} B.~S.~Acharya, S.~Gukov, ``M-theory and singularities of exceptional holonomy manifolds,'' \textit{Phys.\ Rep.} \textbf{392}, 121 (2004).

\bibitem{12} L.~Foscolo et al., ``Complete non-compact $G_2$-manifolds from asymptotically conical Calabi--Yau 3-folds,'' \textit{Duke Math.\ J.} \textbf{170}, 3 (2021).

\bibitem{13} A.~Kovalev, ``Coassociative K3 fibrations of compact $G_2$-manifolds,'' arXiv:math/0511150 (2005).

\bibitem{14} M.~Haskins, J.~Nordstr\"om, ``Extra-twisted connected sum $G_2$-manifolds,'' arXiv:1809.09083 (2022); \textit{Ann.\ Glob.\ Anal.\ Geom.} \textbf{64}, art.\ 2 (2023).

\bibitem{15} G.~Bera, ``Associative submanifolds in twisted connected sum $G_2$-manifolds,'' arXiv:2209.00156 (2022).

\bibitem{16} D.~Crowley, S.~Goette, J.~Nordstr\"om, ``An analytic invariant of $G_2$ manifolds,'' \textit{Invent.\ Math.} \textbf{239}(3), 865--907 (2025).

\bibitem{17} T.~Dray, C.~A.~Manogue, \textit{The Geometry of the Octonions}, World Scientific (2015).

\bibitem{18} J.~F.~Adams, \textit{Lectures on Exceptional Lie Groups}, University of Chicago Press (1996).

\bibitem{19} D.~J.~Gross et al., ``Heterotic string theory,'' \textit{Nucl.\ Phys.\ B} \textbf{256}, 253 (1985).

\bibitem{20} D.~D.~Joyce, \textit{Compact Manifolds with Special Holonomy}, Oxford University Press (2000).

\bibitem{21} A.~Kovalev, ``Twisted connected sums and special Riemannian holonomy,'' \textit{J.\ Reine Angew.\ Math.} \textbf{565}, 125 (2003).

\bibitem{22} A.~Corti et al., ``$G_2$-manifolds and associative submanifolds via semi-Fano 3-folds,'' \textit{Duke Math.\ J.} \textbf{164}, 1971 (2015).

\bibitem{23} R.~Harvey, H.~B.~Lawson, ``Calibrated geometries,'' \textit{Acta Math.} \textbf{148}, 47 (1982).

\bibitem{24} B.~S.~Acharya, ``M-theory, $G_2$-manifolds and four-dimensional physics,'' \textit{Class.\ Quant.\ Grav.} \textbf{19}, 5619 (2002).

\bibitem{25} B.~S.~Acharya, E.~Witten, ``Chiral fermions from manifolds of $G_2$ holonomy,'' arXiv:hep-th/0109152 (2001).

\bibitem{26} T2K and NOvA Collaborations, ``Joint analysis of long-baseline neutrino oscillations,'' \textit{Nature} \textbf{646}, 818--824 (2025), \href{https://doi.org/10.1038/s41586-025-09599-3}{doi:10.1038/s41586-025-09599-3}.

\bibitem{27} I.~Esteban et al.\ (NuFIT), ``NuFit-6.0: three-flavour global analysis of neutrino oscillation experiments,'' \textit{JHEP} \textbf{12} (2024) 216, arXiv:2410.05380; NuFIT 6.1 web release, \url{https://www.nu-fit.org} (2025).

\bibitem{28} L.~de Moura, S.~Ullrich, ``The Lean 4 theorem prover and programming language,'' in \textit{CADE 28}, p.\ 625 (2021).

\bibitem{29} mathlib Community, \textit{mathlib4}, \url{https://github.com/leanprover-community/mathlib4}.

\bibitem{30} DUNE Collaboration, ``Long-baseline neutrino facility (LBNF) and DUNE conceptual design report,'' FERMILAB-TM-2696 (2020).

\bibitem{31} DUNE Collaboration, ``Deep Underground Neutrino Experiment (DUNE) far detector technical design report,'' arXiv:2103.04797 (2021).

\bibitem{32} E.~Witten, ``Strong coupling expansion of Calabi-Yau compactification,'' \textit{Nucl.\ Phys.\ B} \textbf{471}, 135 (1996).

\bibitem{33} B.~S.~Acharya et al., ``M-theory solution to the hierarchy problem,'' \textit{Phys.\ Rev.\ D} \textbf{76}, 126010 (2007).

\bibitem{34} M.~Atiyah, E.~Witten, ``M-theory dynamics on a manifold of $G_2$-holonomy,'' \textit{Adv.\ Theor.\ Math.\ Phys.} \textbf{6}, 1 (2002).

\bibitem{35} G.~Kane, \textit{String Theory and the Real World} (2017).

\bibitem{36} J.~Distler, S.~Garibaldi, ``There is no `Theory of Everything' inside $E_8$,'' \textit{Commun.\ Math.\ Phys.} \textbf{298}, 419 (2010).

\bibitem{37} J.~C.~Baez, ``Octonions and the Standard Model,'' \url{https://math.ucr.edu/home/baez/standard/} (2020--2025).

\bibitem{38} Springer Nature, ``Artificial intelligence (AI) policy,'' \url{https://www.springernature.com/gp/policies} (2024).

\bibitem{39} J.~A.~Wheeler, ``Information, physics, quantum: the search for links,'' in \textit{Complexity, Entropy, and the Physics of Information}, W.~H.~Zurek (ed.), Addison-Wesley (1990), pp.~3--28.

\bibitem{40} J.~Worrall, ``Structural realism: the best of both worlds?,'' \textit{Dialectica} \textbf{43}, 99--124 (1989).

\bibitem{41} J.~Ladyman, D.~Ross, \textit{Every Thing Must Go: Metaphysics Naturalized}, Oxford University Press (2007).

\bibitem{42} R.~Pin\v{c}\'ak, A.~Pigazzini, M.~Pudl\'ak, E.~Barto\v{s}, ``Geometric origin of a stable black hole remnant from torsion in $G_2$-manifold geometry,'' \textit{Gen.\ Rel.\ Grav.} \textbf{58}, 29 (2026), \href{https://doi.org/10.1007/s10714-026-03528-z}{doi:10.1007/s10714-026-03528-z}.

\bibitem{43} D.~Joyce, S.~Karigiannis, ``A new construction of compact torsion-free $G_2$-manifolds by gluing families of Eguchi--Hanson spaces,'' \textit{J.\ Differential Geom.} (2021), arXiv:1707.09325 (2017).

\bibitem{44} S.~Mukai, ``Finite groups of automorphisms of K3 surfaces and the Mathieu group,'' \textit{Invent.\ Math.} \textbf{94}, 183--221 (1988).

\bibitem{45} A.~Garbagnati, A.~Sarti, ``Symplectic automorphisms of prime order on K3 surfaces,'' \textit{J.\ Algebra} \textbf{318}, 323--350 (2007), arXiv:math/0603742.

\bibitem{46} LEP SUSY Working Group (ALEPH, DELPHI, L3, OPAL), ``Combined LEP chargino results, up to 208 GeV for large $m_0$,'' LEPSUSYWG/01-03.1 (2001), \url{https://lepsusy.web.cern.ch/lepsusy/www/inos_moriond01/charginos_pub.html}.

\bibitem{47} Fermilab, ``DUNE first-beam target of 2031: LBNF/DUNE-US milestone communication'' (2026-05-07), \url{https://news.fnal.gov/}.

\bibitem{48} J.~Chen, H.~Hong, ``Intermediate curvature and splitting theorem,'' arXiv:2604.26529 (2026).

\bibitem{49} M.~Tristram et al., ``Cosmological parameters derived from the final Planck data release (PR4),'' \textit{Astron.\ Astrophys.} \textbf{682}, A37 (2024).

\bibitem{50} E.~Tiesinga, P.~J.~Mohr, D.~B.~Newell, B.~N.~Taylor, ``CODATA recommended values of the fundamental physical constants: 2022,'' \textit{Rev.\ Mod.\ Phys.} \textbf{97}, 025002 (2025), \url{https://physics.nist.gov/cuu/Constants}.

\bibitem{51} V.~V.~Nikulin, ``Integer symmetric bilinear forms and some of their geometric applications,'' \textit{Izv.\ Akad.\ Nauk SSSR Ser.\ Mat.} \textbf{43} (1979); English transl.\ \textit{Math.\ USSR Izv.} \textbf{14} (1980).

\end{thebibliography}
 where the references section should appear.
% Sequential numbering [1]..[50]. Round-1 + round-2 IA-review updates applied
% (2026-07-07): [13] Kovalev 2005, [26] T2K/NOvA Nature 646,
% [27] NuFit 6.0 + 6.1, [45] Garbagnati-Sarti 2007, [46] LEPSUSYWG,
% [47] Fermilab DUNE 2031, [48] Chen-Hong (re-added, cited in S1), [49] Planck PR4, [50] CODATA.

\begin{thebibliography}{99}

\bibitem{1} Particle Data Group, ``Review of Particle Physics,'' \textit{Phys.\ Rev.\ D} \textbf{110}, 030001 (2024).

\bibitem{2} S.~Weinberg, ``Implications of dynamical symmetry breaking,'' \textit{Phys.\ Rev.\ D} \textbf{13}, 974 (1976).

\bibitem{3} Planck Collaboration, ``Planck 2018 results.\ VI.\ Cosmological parameters,'' \textit{Astron.\ Astrophys.} \textbf{641}, A6 (2020).

\bibitem{4} A.~G.~Riess et al., ``A comprehensive measurement of the local value of the Hubble constant,'' \textit{ApJL} \textbf{934}, L7 (2022).

\bibitem{5} C.~D.~Froggatt, H.~B.~Nielsen, ``Hierarchy of quark masses, Cabibbo angles and CP violation,'' \textit{Nucl.\ Phys.\ B} \textbf{147}, 277 (1979).

\bibitem{6} Y.~Koide, ``Fermion-boson two-body model of quarks and leptons and Cabibbo mixing,'' \textit{Lett.\ Nuovo Cim.} \textbf{34}, 201 (1982).

\bibitem{7} C.~Furey, \textit{Standard Model Physics from an Algebra?}, PhD thesis, University of Waterloo (2015).

\bibitem{8} N.~Furey, M.~J.~Hughes, ``One generation of standard model Weyl representations as a single copy of $\mathbb{R} \otimes \mathbb{C} \otimes \mathbb{H} \otimes \mathbb{O}$,'' \textit{Phys.\ Lett.\ B} \textbf{831}, 137186 (2022).

\bibitem{9} R.~Wilson, ``$E_8$ and Standard Model plus gravity,'' arXiv:2404.18938 [hep-th] (2024).

\bibitem{10} T.~P.~Singh et al., ``An $E_8 \otimes E_8$ unification of the Standard Model with pre-gravitation,'' arXiv:2206.06911v3 (2024).

\bibitem{11} B.~S.~Acharya, S.~Gukov, ``M-theory and singularities of exceptional holonomy manifolds,'' \textit{Phys.\ Rep.} \textbf{392}, 121 (2004).

\bibitem{12} L.~Foscolo et al., ``Complete non-compact $G_2$-manifolds from asymptotically conical Calabi--Yau 3-folds,'' \textit{Duke Math.\ J.} \textbf{170}, 3 (2021).

\bibitem{13} A.~Kovalev, ``Coassociative K3 fibrations of compact $G_2$-manifolds,'' arXiv:math/0511150 (2005).

\bibitem{14} M.~Haskins, J.~Nordstr\"om, ``Extra-twisted connected sum $G_2$-manifolds,'' arXiv:1809.09083 (2022); \textit{Ann.\ Glob.\ Anal.\ Geom.} \textbf{64}, art.\ 2 (2023).

\bibitem{15} G.~Bera, ``Associative submanifolds in twisted connected sum $G_2$-manifolds,'' arXiv:2209.00156 (2022).

\bibitem{16} D.~Crowley, S.~Goette, J.~Nordstr\"om, ``An analytic invariant of $G_2$ manifolds,'' \textit{Invent.\ Math.} \textbf{239}(3), 865--907 (2025).

\bibitem{17} T.~Dray, C.~A.~Manogue, \textit{The Geometry of the Octonions}, World Scientific (2015).

\bibitem{18} J.~F.~Adams, \textit{Lectures on Exceptional Lie Groups}, University of Chicago Press (1996).

\bibitem{19} D.~J.~Gross et al., ``Heterotic string theory,'' \textit{Nucl.\ Phys.\ B} \textbf{256}, 253 (1985).

\bibitem{20} D.~D.~Joyce, \textit{Compact Manifolds with Special Holonomy}, Oxford University Press (2000).

\bibitem{21} A.~Kovalev, ``Twisted connected sums and special Riemannian holonomy,'' \textit{J.\ Reine Angew.\ Math.} \textbf{565}, 125 (2003).

\bibitem{22} A.~Corti et al., ``$G_2$-manifolds and associative submanifolds via semi-Fano 3-folds,'' \textit{Duke Math.\ J.} \textbf{164}, 1971 (2015).

\bibitem{23} R.~Harvey, H.~B.~Lawson, ``Calibrated geometries,'' \textit{Acta Math.} \textbf{148}, 47 (1982).

\bibitem{24} B.~S.~Acharya, ``M-theory, $G_2$-manifolds and four-dimensional physics,'' \textit{Class.\ Quant.\ Grav.} \textbf{19}, 5619 (2002).

\bibitem{25} B.~S.~Acharya, E.~Witten, ``Chiral fermions from manifolds of $G_2$ holonomy,'' arXiv:hep-th/0109152 (2001).

\bibitem{26} T2K and NOvA Collaborations, ``Joint analysis of long-baseline neutrino oscillations,'' \textit{Nature} \textbf{646}, 818--824 (2025), \href{https://doi.org/10.1038/s41586-025-09599-3}{doi:10.1038/s41586-025-09599-3}.

\bibitem{27} I.~Esteban et al.\ (NuFIT), ``NuFit-6.0: three-flavour global analysis of neutrino oscillation experiments,'' \textit{JHEP} \textbf{12} (2024) 216, arXiv:2410.05380; NuFIT 6.1 web release, \url{https://www.nu-fit.org} (2025).

\bibitem{28} L.~de Moura, S.~Ullrich, ``The Lean 4 theorem prover and programming language,'' in \textit{CADE 28}, p.\ 625 (2021).

\bibitem{29} mathlib Community, \textit{mathlib4}, \url{https://github.com/leanprover-community/mathlib4}.

\bibitem{30} DUNE Collaboration, ``Long-baseline neutrino facility (LBNF) and DUNE conceptual design report,'' FERMILAB-TM-2696 (2020).

\bibitem{31} DUNE Collaboration, ``Deep Underground Neutrino Experiment (DUNE) far detector technical design report,'' arXiv:2103.04797 (2021).

\bibitem{32} E.~Witten, ``Strong coupling expansion of Calabi-Yau compactification,'' \textit{Nucl.\ Phys.\ B} \textbf{471}, 135 (1996).

\bibitem{33} B.~S.~Acharya et al., ``M-theory solution to the hierarchy problem,'' \textit{Phys.\ Rev.\ D} \textbf{76}, 126010 (2007).

\bibitem{34} M.~Atiyah, E.~Witten, ``M-theory dynamics on a manifold of $G_2$-holonomy,'' \textit{Adv.\ Theor.\ Math.\ Phys.} \textbf{6}, 1 (2002).

\bibitem{35} G.~Kane, \textit{String Theory and the Real World} (2017).

\bibitem{36} J.~Distler, S.~Garibaldi, ``There is no `Theory of Everything' inside $E_8$,'' \textit{Commun.\ Math.\ Phys.} \textbf{298}, 419 (2010).

\bibitem{37} J.~C.~Baez, ``Octonions and the Standard Model,'' \url{https://math.ucr.edu/home/baez/standard/} (2020--2025).

\bibitem{38} Springer Nature, ``Artificial intelligence (AI) policy,'' \url{https://www.springernature.com/gp/policies} (2024).

\bibitem{39} J.~A.~Wheeler, ``Information, physics, quantum: the search for links,'' in \textit{Complexity, Entropy, and the Physics of Information}, W.~H.~Zurek (ed.), Addison-Wesley (1990), pp.~3--28.

\bibitem{40} J.~Worrall, ``Structural realism: the best of both worlds?,'' \textit{Dialectica} \textbf{43}, 99--124 (1989).

\bibitem{41} J.~Ladyman, D.~Ross, \textit{Every Thing Must Go: Metaphysics Naturalized}, Oxford University Press (2007).

\bibitem{42} R.~Pin\v{c}\'ak, A.~Pigazzini, M.~Pudl\'ak, E.~Barto\v{s}, ``Geometric origin of a stable black hole remnant from torsion in $G_2$-manifold geometry,'' \textit{Gen.\ Rel.\ Grav.} \textbf{58}, 29 (2026), \href{https://doi.org/10.1007/s10714-026-03528-z}{doi:10.1007/s10714-026-03528-z}.

\bibitem{43} D.~Joyce, S.~Karigiannis, ``A new construction of compact torsion-free $G_2$-manifolds by gluing families of Eguchi--Hanson spaces,'' \textit{J.\ Differential Geom.} (2021), arXiv:1707.09325 (2017).

\bibitem{44} S.~Mukai, ``Finite groups of automorphisms of K3 surfaces and the Mathieu group,'' \textit{Invent.\ Math.} \textbf{94}, 183--221 (1988).

\bibitem{45} A.~Garbagnati, A.~Sarti, ``Symplectic automorphisms of prime order on K3 surfaces,'' \textit{J.\ Algebra} \textbf{318}, 323--350 (2007), arXiv:math/0603742.

\bibitem{46} LEP SUSY Working Group (ALEPH, DELPHI, L3, OPAL), ``Combined LEP chargino results, up to 208 GeV for large $m_0$,'' LEPSUSYWG/01-03.1 (2001), \url{https://lepsusy.web.cern.ch/lepsusy/www/inos_moriond01/charginos_pub.html}.

\bibitem{47} Fermilab, ``DUNE first-beam target of 2031: LBNF/DUNE-US milestone communication'' (2026-05-07), \url{https://news.fnal.gov/}.

\bibitem{48} J.~Chen, H.~Hong, ``Intermediate curvature and splitting theorem,'' arXiv:2604.26529 (2026).

\bibitem{49} M.~Tristram et al., ``Cosmological parameters derived from the final Planck data release (PR4),'' \textit{Astron.\ Astrophys.} \textbf{682}, A37 (2024).

\bibitem{50} E.~Tiesinga, P.~J.~Mohr, D.~B.~Newell, B.~N.~Taylor, ``CODATA recommended values of the fundamental physical constants: 2022,'' \textit{Rev.\ Mod.\ Phys.} \textbf{97}, 025002 (2025), \url{https://physics.nist.gov/cuu/Constants}.

\bibitem{51} V.~V.~Nikulin, ``Integer symmetric bilinear forms and some of their geometric applications,'' \textit{Izv.\ Akad.\ Nauk SSSR Ser.\ Mat.} \textbf{43} (1979); English transl.\ \textit{Math.\ USSR Izv.} \textbf{14} (1980).

\end{thebibliography}
 where the references section should appear.
% Sequential numbering [1]..[50]. Round-1 + round-2 IA-review updates applied
% (2026-07-07): [13] Kovalev 2005, [26] T2K/NOvA Nature 646,
% [27] NuFit 6.0 + 6.1, [45] Garbagnati-Sarti 2007, [46] LEPSUSYWG,
% [47] Fermilab DUNE 2031, [48] Chen-Hong (re-added, cited in S1), [49] Planck PR4, [50] CODATA.

\begin{thebibliography}{99}

\bibitem{1} Particle Data Group, ``Review of Particle Physics,'' \textit{Phys.\ Rev.\ D} \textbf{110}, 030001 (2024).

\bibitem{2} S.~Weinberg, ``Implications of dynamical symmetry breaking,'' \textit{Phys.\ Rev.\ D} \textbf{13}, 974 (1976).

\bibitem{3} Planck Collaboration, ``Planck 2018 results.\ VI.\ Cosmological parameters,'' \textit{Astron.\ Astrophys.} \textbf{641}, A6 (2020).

\bibitem{4} A.~G.~Riess et al., ``A comprehensive measurement of the local value of the Hubble constant,'' \textit{ApJL} \textbf{934}, L7 (2022).

\bibitem{5} C.~D.~Froggatt, H.~B.~Nielsen, ``Hierarchy of quark masses, Cabibbo angles and CP violation,'' \textit{Nucl.\ Phys.\ B} \textbf{147}, 277 (1979).

\bibitem{6} Y.~Koide, ``Fermion-boson two-body model of quarks and leptons and Cabibbo mixing,'' \textit{Lett.\ Nuovo Cim.} \textbf{34}, 201 (1982).

\bibitem{7} C.~Furey, \textit{Standard Model Physics from an Algebra?}, PhD thesis, University of Waterloo (2015).

\bibitem{8} N.~Furey, M.~J.~Hughes, ``One generation of standard model Weyl representations as a single copy of $\mathbb{R} \otimes \mathbb{C} \otimes \mathbb{H} \otimes \mathbb{O}$,'' \textit{Phys.\ Lett.\ B} \textbf{831}, 137186 (2022).

\bibitem{9} R.~Wilson, ``$E_8$ and Standard Model plus gravity,'' arXiv:2404.18938 [hep-th] (2024).

\bibitem{10} T.~P.~Singh et al., ``An $E_8 \otimes E_8$ unification of the Standard Model with pre-gravitation,'' arXiv:2206.06911v3 (2024).

\bibitem{11} B.~S.~Acharya, S.~Gukov, ``M-theory and singularities of exceptional holonomy manifolds,'' \textit{Phys.\ Rep.} \textbf{392}, 121 (2004).

\bibitem{12} L.~Foscolo et al., ``Complete non-compact $G_2$-manifolds from asymptotically conical Calabi--Yau 3-folds,'' \textit{Duke Math.\ J.} \textbf{170}, 3 (2021).

\bibitem{13} A.~Kovalev, ``Coassociative K3 fibrations of compact $G_2$-manifolds,'' arXiv:math/0511150 (2005).

\bibitem{14} M.~Haskins, J.~Nordstr\"om, ``Extra-twisted connected sum $G_2$-manifolds,'' arXiv:1809.09083 (2022); \textit{Ann.\ Glob.\ Anal.\ Geom.} \textbf{64}, art.\ 2 (2023).

\bibitem{15} G.~Bera, ``Associative submanifolds in twisted connected sum $G_2$-manifolds,'' arXiv:2209.00156 (2022).

\bibitem{16} D.~Crowley, S.~Goette, J.~Nordstr\"om, ``An analytic invariant of $G_2$ manifolds,'' \textit{Invent.\ Math.} \textbf{239}(3), 865--907 (2025).

\bibitem{17} T.~Dray, C.~A.~Manogue, \textit{The Geometry of the Octonions}, World Scientific (2015).

\bibitem{18} J.~F.~Adams, \textit{Lectures on Exceptional Lie Groups}, University of Chicago Press (1996).

\bibitem{19} D.~J.~Gross et al., ``Heterotic string theory,'' \textit{Nucl.\ Phys.\ B} \textbf{256}, 253 (1985).

\bibitem{20} D.~D.~Joyce, \textit{Compact Manifolds with Special Holonomy}, Oxford University Press (2000).

\bibitem{21} A.~Kovalev, ``Twisted connected sums and special Riemannian holonomy,'' \textit{J.\ Reine Angew.\ Math.} \textbf{565}, 125 (2003).

\bibitem{22} A.~Corti et al., ``$G_2$-manifolds and associative submanifolds via semi-Fano 3-folds,'' \textit{Duke Math.\ J.} \textbf{164}, 1971 (2015).

\bibitem{23} R.~Harvey, H.~B.~Lawson, ``Calibrated geometries,'' \textit{Acta Math.} \textbf{148}, 47 (1982).

\bibitem{24} B.~S.~Acharya, ``M-theory, $G_2$-manifolds and four-dimensional physics,'' \textit{Class.\ Quant.\ Grav.} \textbf{19}, 5619 (2002).

\bibitem{25} B.~S.~Acharya, E.~Witten, ``Chiral fermions from manifolds of $G_2$ holonomy,'' arXiv:hep-th/0109152 (2001).

\bibitem{26} T2K and NOvA Collaborations, ``Joint analysis of long-baseline neutrino oscillations,'' \textit{Nature} \textbf{646}, 818--824 (2025), \href{https://doi.org/10.1038/s41586-025-09599-3}{doi:10.1038/s41586-025-09599-3}.

\bibitem{27} I.~Esteban et al.\ (NuFIT), ``NuFit-6.0: three-flavour global analysis of neutrino oscillation experiments,'' \textit{JHEP} \textbf{12} (2024) 216, arXiv:2410.05380; NuFIT 6.1 web release, \url{https://www.nu-fit.org} (2025).

\bibitem{28} L.~de Moura, S.~Ullrich, ``The Lean 4 theorem prover and programming language,'' in \textit{CADE 28}, p.\ 625 (2021).

\bibitem{29} mathlib Community, \textit{mathlib4}, \url{https://github.com/leanprover-community/mathlib4}.

\bibitem{30} DUNE Collaboration, ``Long-baseline neutrino facility (LBNF) and DUNE conceptual design report,'' FERMILAB-TM-2696 (2020).

\bibitem{31} DUNE Collaboration, ``Deep Underground Neutrino Experiment (DUNE) far detector technical design report,'' arXiv:2103.04797 (2021).

\bibitem{32} E.~Witten, ``Strong coupling expansion of Calabi-Yau compactification,'' \textit{Nucl.\ Phys.\ B} \textbf{471}, 135 (1996).

\bibitem{33} B.~S.~Acharya et al., ``M-theory solution to the hierarchy problem,'' \textit{Phys.\ Rev.\ D} \textbf{76}, 126010 (2007).

\bibitem{34} M.~Atiyah, E.~Witten, ``M-theory dynamics on a manifold of $G_2$-holonomy,'' \textit{Adv.\ Theor.\ Math.\ Phys.} \textbf{6}, 1 (2002).

\bibitem{35} G.~Kane, \textit{String Theory and the Real World} (2017).

\bibitem{36} J.~Distler, S.~Garibaldi, ``There is no `Theory of Everything' inside $E_8$,'' \textit{Commun.\ Math.\ Phys.} \textbf{298}, 419 (2010).

\bibitem{37} J.~C.~Baez, ``Octonions and the Standard Model,'' \url{https://math.ucr.edu/home/baez/standard/} (2020--2025).

\bibitem{38} Springer Nature, ``Artificial intelligence (AI) policy,'' \url{https://www.springernature.com/gp/policies} (2024).

\bibitem{39} J.~A.~Wheeler, ``Information, physics, quantum: the search for links,'' in \textit{Complexity, Entropy, and the Physics of Information}, W.~H.~Zurek (ed.), Addison-Wesley (1990), pp.~3--28.

\bibitem{40} J.~Worrall, ``Structural realism: the best of both worlds?,'' \textit{Dialectica} \textbf{43}, 99--124 (1989).

\bibitem{41} J.~Ladyman, D.~Ross, \textit{Every Thing Must Go: Metaphysics Naturalized}, Oxford University Press (2007).

\bibitem{42} R.~Pin\v{c}\'ak, A.~Pigazzini, M.~Pudl\'ak, E.~Barto\v{s}, ``Geometric origin of a stable black hole remnant from torsion in $G_2$-manifold geometry,'' \textit{Gen.\ Rel.\ Grav.} \textbf{58}, 29 (2026), \href{https://doi.org/10.1007/s10714-026-03528-z}{doi:10.1007/s10714-026-03528-z}.

\bibitem{43} D.~Joyce, S.~Karigiannis, ``A new construction of compact torsion-free $G_2$-manifolds by gluing families of Eguchi--Hanson spaces,'' \textit{J.\ Differential Geom.} (2021), arXiv:1707.09325 (2017).

\bibitem{44} S.~Mukai, ``Finite groups of automorphisms of K3 surfaces and the Mathieu group,'' \textit{Invent.\ Math.} \textbf{94}, 183--221 (1988).

\bibitem{45} A.~Garbagnati, A.~Sarti, ``Symplectic automorphisms of prime order on K3 surfaces,'' \textit{J.\ Algebra} \textbf{318}, 323--350 (2007), arXiv:math/0603742.

\bibitem{46} LEP SUSY Working Group (ALEPH, DELPHI, L3, OPAL), ``Combined LEP chargino results, up to 208 GeV for large $m_0$,'' LEPSUSYWG/01-03.1 (2001), \url{https://lepsusy.web.cern.ch/lepsusy/www/inos_moriond01/charginos_pub.html}.

\bibitem{47} Fermilab, ``DUNE first-beam target of 2031: LBNF/DUNE-US milestone communication'' (2026-05-07), \url{https://news.fnal.gov/}.

\bibitem{48} J.~Chen, H.~Hong, ``Intermediate curvature and splitting theorem,'' arXiv:2604.26529 (2026).

\bibitem{49} M.~Tristram et al., ``Cosmological parameters derived from the final Planck data release (PR4),'' \textit{Astron.\ Astrophys.} \textbf{682}, A37 (2024).

\bibitem{50} E.~Tiesinga, P.~J.~Mohr, D.~B.~Newell, B.~N.~Taylor, ``CODATA recommended values of the fundamental physical constants: 2022,'' \textit{Rev.\ Mod.\ Phys.} \textbf{97}, 025002 (2025), \url{https://physics.nist.gov/cuu/Constants}.

\bibitem{51} V.~V.~Nikulin, ``Integer symmetric bilinear forms and some of their geometric applications,'' \textit{Izv.\ Akad.\ Nauk SSSR Ser.\ Mat.} \textbf{43} (1979); English transl.\ \textit{Math.\ USSR Izv.} \textbf{14} (1980).

\end{thebibliography}


\end{document}
